% =====================================================================
%  DIRECCIÓN FINANCIERA (UGR) · Relación de problemas resuelta · P1–P25
%  Fuente LaTeX del PDF publicado en
%  public/ade/direccion-financiera/relacion-problemas-resueltos.pdf
%  Compilar:  ./compilar.sh   (pdflatex x2)
% =====================================================================
\documentclass[10pt,a4paper,twoside]{article}

\usepackage[a4paper,top=25mm,bottom=22mm,inner=21mm,outer=19mm,
            headheight=19pt,headsep=6.3mm,footskip=10mm]{geometry}
\usepackage[T1]{fontenc}
\usepackage[utf8]{inputenc}
\usepackage[spanish,es-noshorthands,es-nodecimaldot]{babel}
\AtBeginDocument{\spanishplainpercent}
\usepackage{sourceserifpro}
\usepackage[scale=0.96]{sourcesanspro}
\usepackage{amsmath,amssymb}
\usepackage{icomma}
\usepackage[final]{microtype}
\usepackage{xcolor}
\usepackage{tikz}
\usetikzlibrary{arrows.meta,calc,positioning,decorations.pathreplacing}
\usepackage{pgfplots}
\pgfplotsset{compat=1.18}
\usepackage[most]{tcolorbox}
\usepackage{booktabs,tabularx,colortbl,array,multirow}
\usepackage{enumitem}
\usepackage{titlesec}
\usepackage{fancyhdr}
\usepackage[breakall]{truncate}
\usepackage{fontawesome5}
\usepackage{needspace}
\usepackage{xspace}
\raggedbottom
\usepackage[hidelinks,bookmarksnumbered]{hyperref}

% ---------------------------------------------------------------------
%  Identidad corporativa (idéntica a los apuntes)
% ---------------------------------------------------------------------
\definecolor{navy}{HTML}{0B2545}
\definecolor{steel}{HTML}{13315C}
\definecolor{slate}{HTML}{5B6B7F}
\definecolor{mist}{HTML}{EEF2F6}
\definecolor{rule}{HTML}{C9D3DE}
\definecolor{gold}{HTML}{B88A1B}
\definecolor{goldbg}{HTML}{FBF5E6}
\definecolor{alert}{HTML}{A4243B}
\definecolor{alertbg}{HTML}{FBEEF0}
\definecolor{okgreen}{HTML}{2E7D5B}
\definecolor{okbg}{HTML}{EDF6F1}

\setlength{\parindent}{0pt}
\setlength{\parskip}{4pt plus 1pt}
\linespread{1.06}
\setlist{itemsep=2pt,topsep=3pt,leftmargin=1.3em}
\setlist[itemize]{label={\color{steel}\scriptsize$\blacksquare$}}
\setlist[enumerate]{label={\sffamily\bfseries\color{steel}\arabic*.}}

% Notación oficial del Departamento
\newcommand{\di}{)_{di}}
\newcommand{\ai}{)_{ai}}
\newcommand{\an}[2]{a_{\overline{#1}|#2}}
\newcommand{\eur}{\,\text{€}\xspace}
\newcommand{\paso}[1]{\par\smallskip{\sffamily\bfseries\color{steel}#1}\par\nopagebreak}

% ---------------------------------------------------------------------
%  Títulos
% ---------------------------------------------------------------------
\titleformat{\section}[block]
  {\sffamily\bfseries\Large\color{navy}}
  {}{0pt}{}[\vspace{2pt}{\color{gold}\titlerule[1.2pt]}]
\titlespacing*{\section}{0pt}{16pt}{8pt}

\newcommand{\bloque}[3]{%
  \par\needspace{0.42\textheight}\vspace{6pt}
  \phantomsection
  \addcontentsline{toc}{part}{#1 · #2}%
  \markboth{#1 · #2}{}%
  \begin{tcolorbox}[enhanced,colback=navy,colframe=navy,arc=0pt,
      left=12pt,right=12pt,top=14pt,bottom=14pt,boxrule=0pt,
      borderline north={2.5pt}{0pt}{gold}]
    {\sffamily\color{gold}\bfseries\small #1}\par\vspace{3pt}
    {\sffamily\color{white}\bfseries\LARGE #2}\par\vspace{6pt}
    {\sffamily\color{white!80!navy}\small #3}
  \end{tcolorbox}\vspace{2pt}%
}

% Cabecera de problema: P# + título + etiquetas
\newcommand{\problema}[3]{%
  \needspace{16\baselineskip}%
  \phantomsection
  \addcontentsline{toc}{section}{P#1 · #2}%
  \section*{\textcolor{gold}{P#1}\enspace#2}%
  \vspace{-6pt}{\sffamily\footnotesize\color{slate}#3}\par\vspace{2pt}%
}

% ---------------------------------------------------------------------
%  Cajas
% ---------------------------------------------------------------------
\tcbset{base/.style={enhanced,breakable,arc=1.5pt,boxrule=0pt,
  left=8pt,right=8pt,top=5pt,bottom=5pt,before skip=7pt,after skip=7pt,
  fonttitle=\sffamily\bfseries\small,coltitle=white}}

\newtcolorbox{enunciado}[1]{base,colback=mist,colframe=steel,
  borderline west={3pt}{0pt}{steel},
  title={\faFile*\ \ Enunciado oficial · Problema #1},colbacktitle=steel}
\newtcolorbox{solucion}[1][Resultado]{base,colback=okbg,colframe=okgreen,
  borderline west={3pt}{0pt}{okgreen},
  title={\faCheckCircle\ \ #1},colbacktitle=okgreen}
\newtcolorbox{clave}[1][Idea clave]{base,colback=goldbg,colframe=gold,
  borderline west={3pt}{0pt}{gold},title={\faLightbulb\ \ #1},
  colbacktitle=gold}
\newtcolorbox{trampa}[1][Ojo]{base,colback=alertbg,colframe=alert,
  borderline west={3pt}{0pt}{alert},title={\faExclamationTriangle\ \ #1},
  colbacktitle=alert}

% Opciones de test
\newenvironment{opciones}{\begin{enumerate}[label={\sffamily\bfseries\color{steel}\alph*)},
  leftmargin=1.6em,itemsep=1pt,topsep=2pt]}{\end{enumerate}}

% Tablas
\newcolumntype{L}{>{\raggedright\arraybackslash}X}
\newcolumntype{R}{>{\raggedleft\arraybackslash}X}
\newcolumntype{Q}{>{\raggedleft\arraybackslash}p{1.95cm}}
\newcolumntype{N}[1]{>{\raggedleft\arraybackslash}p{#1}}
\newcommand{\thead}[1]{\leavevmode\sffamily\bfseries\color{white}#1}
\newcommand{\headrow}{\rowcolor{steel}}
\newcommand{\totrow}{\rowcolor{goldbg}}
\renewcommand{\arraystretch}{1.18}

% Línea temporal de flujos: \flujos[ancho por periodo]{n}{t/s/texto,...}
%   s = p (cobro, flecha hacia arriba) · s = n (pago, flecha hacia abajo)
\newcommand{\flujos}[3][1.9]{%
\begin{tikzpicture}[x=#1cm,font=\sffamily\scriptsize]
  \draw[slate,line width=0.7pt,-{Stealth[length=2mm]}] (-0.35,0) -- (#2+0.55,0);
  \foreach \t in {0,...,#2}{\draw[slate] (\t,0.07) -- (\t,-0.07);
    \node[slate,fill=white,inner sep=1pt] at (\t,-0.28) {$t=\t$};}
  \foreach \t/\s/\v in {#3}{%
    \if\s p
      \draw[okgreen,line width=1pt,-{Stealth[length=2mm]}] (\t,0.12) -- (\t,0.85)
        node[above,align=center,text=okgreen!70!black]{\v};
    \else
      \draw[alert,line width=1pt,-{Stealth[length=2mm]}] (\t,-0.45) -- (\t,-1.15)
        node[below,align=center,text=alert]{\v};
    \fi}
\end{tikzpicture}}

% ---------------------------------------------------------------------
%  Cabeceras y pies
% ---------------------------------------------------------------------
\pagestyle{fancy}
\fancyhf{}
\renewcommand{\headrulewidth}{0.4pt}
\renewcommand{\headrule}{\hbox to\headwidth{\color{rule}\leaders\hrule height \headrulewidth\hfill}}
\fancyhead[LE]{\sffamily\footnotesize\bfseries\color{navy}\thepage\hspace{1.2em}%
  \mdseries\color{slate}\truncate{0.78\textwidth}{Dirección Financiera · Relación de problemas resuelta P1–P25}}
\fancyhead[RO]{\sffamily\footnotesize\color{slate}\truncate{0.78\textwidth}{\leftmark}%
  \hspace{1.2em}\bfseries\color{navy}\thepage}
\fancyfoot[C]{\sffamily\scriptsize\color{slate!80}repasaYA · Grado en ADE · Universidad de Granada}
\fancypagestyle{plain}{\fancyhf{}\renewcommand{\headrulewidth}{0pt}}

\hypersetup{pdftitle={Dirección Financiera · Relación de problemas resuelta (P1–P25)},
  pdfauthor={repasaYA},pdfsubject={Primer parcial · UGR}}

\makeatletter
\renewcommand{\l@part}[2]{\vspace{5pt}\noindent
  {\sffamily\bfseries\color{navy}#1\hfill\color{steel}#2}\par}
\renewcommand{\l@section}{\@dottedtocline{1}{1.2em}{2.4em}}
\setcounter{tocdepth}{1}
\makeatother

% =====================================================================
\begin{document}
% =====================================================================

% ---------------------------------------------------------------------
%  PORTADA
% ---------------------------------------------------------------------
\begin{titlepage}
\thispagestyle{empty}
\begin{tikzpicture}[remember picture,overlay]
  \fill[navy] (current page.north west) rectangle ([yshift=-128mm]current page.north east);
  \fill[gold] ([yshift=-128mm]current page.north west) rectangle ([yshift=-130mm]current page.north east);
  \node[anchor=north west,text width=150mm,align=left] at ([xshift=24mm,yshift=-30mm]current page.north west)
   {\sffamily\color{gold}\bfseries\large DIRECCIÓN FINANCIERA\\[2pt]
    \color{white!75!navy}\mdseries\normalsize Grado en ADE · Facultad de CC.~Económicas y Empresariales · UGR};
  \node[anchor=north west,text width=160mm,align=left] at ([xshift=24mm,yshift=-62mm]current page.north west)
   {\sffamily\color{white}\bfseries\fontsize{34}{38}\selectfont Relación de problemas\\[4pt]
    \fontsize{17}{21}\selectfont\mdseries Resuelta · Problemas 1 a 25 · Temas 1 a 4};
  \node[anchor=north west,text width=150mm,align=left] at ([xshift=24mm,yshift=-104mm]current page.north west)
   {\sffamily\color{white!80!navy}\small Enunciados oficiales del curso 2025/26, resolución paso a paso con la\\
    notación del Departamento, gráficos y notas sobre los puntos conflictivos.};
\end{tikzpicture}

\vspace*{118mm}

\begin{minipage}[t]{0.47\textwidth}
{\sffamily\bfseries\color{navy}Contenido}\par\vspace{3pt}
{\color{gold}\rule{\linewidth}{1pt}}\par\vspace{2pt}
\small
\begin{tabularx}{\linewidth}{@{}>{\sffamily\bfseries\color{steel}}l L@{}}
P1--P3   & Fisher: inversión real y mercado de capitales\\
P4--P11  & Acciones: Gordon, multietapa y VAOC\\
P12      & Criterios VAN y TIR\\
P13--P22 & Flujos de caja, VAN e incrementales\\
P23--P25 & CAE y demanda estacional
\end{tabularx}
\end{minipage}\hfill
\begin{minipage}[t]{0.47\textwidth}
{\sffamily\bfseries\color{navy}Cómo usar este documento}\par\vspace{3pt}
{\color{gold}\rule{\linewidth}{1pt}}\par\vspace{2pt}
\small
\begin{itemize}[leftmargin=1.1em]
  \item \textcolor{steel}{\textbf{Azul}}: enunciado oficial. Intenta resolverlo antes de seguir leyendo.
  \item \textcolor{okgreen}{\textbf{Verde}}: resultado final, para comprobar.
  \item \textcolor{gold}{\textbf{Ámbar}}: la idea que hay que llevarse.
  \item \textcolor{alert}{\textbf{Rojo}}: redondeos, erratas y criterios dudosos.
  \item La página 2 trae la plantilla de soluciones.
\end{itemize}
\end{minipage}

\vfill
{\sffamily\scriptsize\color{slate}
Enunciados de la Relación de Problemas de Dirección Financiera 2025/26 (Departamento de Economía Financiera y
Contabilidad, UGR). Resoluciones según las soluciones oficiales de cátedra (P1--P3, P18--P22) y las presentaciones
de los Temas 1--4. Todas las cifras han sido recalculadas.\par}
\end{titlepage}

% ---------------------------------------------------------------------
%  PLANTILLA DE SOLUCIONES + ÍNDICE
% ---------------------------------------------------------------------
\markboth{Plantilla de soluciones}{}
\section*{Plantilla de soluciones}
\addcontentsline{toc}{part}{Plantilla de soluciones}

\begingroup\small
\begin{tabularx}{\textwidth}{@{}>{\leavevmode\sffamily\bfseries\color{navy}}p{0.9cm} L >{\leavevmode\sffamily\bfseries\color{navy}}p{0.9cm} L@{}}
\headrow \thead{P} & \thead{Solución} & \thead{P} & \thead{Solución}\\
P1  & a) & P14 & b) es la falsa\\
\rowcolor{mist}
P2  & b) & P15 & d)\\
P3  & $C_1=6.240$; $\Delta$riqueza $t{=}1$: 2.080; $t{=}0$: 2.000 & P16 & 38.243 bombillas (exacto: 38.244,36)\\
\rowcolor{mist}
P4  & d) & P17 & VAN $=2.317,57$\\
P5  & b) 285,04 & P18 & VAN $=-104,32$\\
\rowcolor{mist}
P6  & b) & P19 & VAN $=112.587,82$\\
P7  & d) & P20 & $\Delta$VAN $=10.388,30$\\
\rowcolor{mist}
P8  & VAOC $=3,19$; $P_0=169,86$; ROE mín.\ 12\,\% & P21 & $\Delta$VAN $=6.346.457,61$ $\Rightarrow$ ampliar\\
P9  & $P_4=749,92$; $P_0=693,25$ & P22 & $\Delta$VAN $=0$; no procede la misma $r$\\
\rowcolor{mist}
P10 & $P_0=115,83$ (VAOC $=2,08$) & P23 & d)\\
P11 & VAOC $=3,64$; $P_0=48,64$; $BPA_1/P_0=9,25\,\%$ & P24 & Horno A (CAE 8.596,99 < 8.996,19)\\
\rowcolor{mist}
P12 & b) & P25 & d) (ahorro 500.000, no 450.000)\\
P13 & c) 732.841,04 & & \\
\end{tabularx}
\endgroup

\begin{clave}[Método común a todos los proyectos (P13--P22)]
\begin{enumerate}
  \item Dibuja la línea temporal y separa tres bloques: $C_0\di$ · $Q_i\di$ · $VR_n\di$.
  \item $C_0\di=$ activo fijo $+$ gastos iniciales$\cdot(1-t_c)$ $+$ FM inicial $[-\,VR_0\di$ del equipo viejo$]$.
  \item $Q_i\di$ renta $=Q_i\ai(1-t_c)+CAF_i\cdot t_c$; \ $Q_i\di$ tesorería $=$ renta $-\Delta FM_i$.
  \item $VR_n\di=VR_n\ai\mp\Delta/\nabla P\cdot t_c$ (plusvalía resta, minusvalía suma) $+$ FM recuperado.
  \item Los intereses y la devolución de préstamos \textbf{no} entran en los flujos: la financiación ya está en $r$.
\end{enumerate}
\end{clave}

\clearpage
\markboth{Índice}{}
{\sffamily\bfseries\large\color{navy}Índice}\par
{\color{gold}\rule{\linewidth}{1.2pt}}\par\vspace{2pt}
\begingroup\small
\makeatletter\@starttoc{toc}\makeatother
\endgroup

% =====================================================================
\clearpage
\bloque{Tema 1}{Separación de Fisher}{Problemas 1 a 3 · inversión en activos reales, mercado de capitales y preferencias de consumo}
% =====================================================================

\problema{1}{Disminución del tipo de interés}{Test · Tema 1 · gráfico de Fisher}

\begin{enunciado}{1}
Una empresa está invirtiendo actualmente en activos reales el importe necesario para maximizar su VAN. Una
disminución de la tasa de interés:
\begin{opciones}
  \item Provocaría un incremento de la inversión inicial necesaria para maximizar el VAN. Un inversor que prefiera
        consumo futuro a consumo actual vería empeorada su situación como consecuencia de la disminución de la tasa de interés.
  \item Provocaría un incremento de la inversión inicial necesaria para maximizar el VAN. Un inversor que prefiera
        consumo actual a consumo futuro vería empeorada su situación como consecuencia de la disminución de la tasa de interés.
  \item Provocaría un incremento de la inversión inicial necesaria para maximizar el VAN. Un inversor que prefiera
        consumo futuro a consumo actual vería mejorada su situación como consecuencia de la disminución de la tasa de interés.
  \item Provocaría una reducción de la inversión inicial necesaria para maximizar el VAN. Un inversor que prefiera
        consumo futuro a consumo actual vería mejorada su situación como consecuencia de la disminución de la tasa de interés.
\end{opciones}
\end{enunciado}

\paso{1. Qué pasa con la inversión}
La empresa invierte hasta el punto en que la curva de inversión en activos reales es tangente a la línea del mercado
de capitales, de pendiente $-(1+r)$. Si $r$ baja, la línea se \textbf{aplana} y la tangencia se desplaza a la
izquierda (de $J$ a $J'$): proyectos que antes rendían menos que $r$ ahora superan el nuevo coste de oportunidad.
\textbf{Aumentan la inversión} ($J'D>JD$) \textbf{y el VAN} ($DK'>DK$).

\begin{center}
\begin{tikzpicture}[font=\sffamily\scriptsize,x=1.1cm,y=0.68cm]
  \draw[-{Stealth[length=2.2mm]},slate] (0,0) -- (9.8,0) node[below left=2pt]{u.m.\ ($t=0$)};
  \draw[-{Stealth[length=2.2mm]},slate] (0,0) -- (0,8.9) node[above right=-1pt]{u.m.\ ($t=1$)};
  % curva de inversión y = 2(6-x)^0,5
  \draw[steel,line width=1.1pt,domain=0.6:6,samples=120,smooth] plot (\x,{2*sqrt(6-\x)});
  % línea con r (pendiente -1,2) tangente en J=5,306
  \draw[gold,line width=1.1pt] (0,8.033) -- (6.694,0);
  % línea con r' (pendiente -0,75) tangente en J'=4,222
  \draw[alert,line width=1.1pt] (0,5.833) -- (7.778,0);
  \fill[navy] (5.306,1.667) circle (1.6pt) node[above right=0pt]{G};
  \fill[navy] (4.222,2.667) circle (1.6pt) node[above right=0pt]{G$'$};
  \draw[dashed,slate] (5.306,1.667) -- (5.306,0);
  \draw[dashed,slate] (4.222,2.667) -- (4.222,0);
  \foreach \x/\l in {4.222/J$'$,5.306/J,6/D,6.694/K,7.778/K$'$}{\node[below] at (\x,0) {\l};}
  \node[left] at (0,8.033) {M};
  \node[left] at (0,5.833) {M$'$};
  \node[gold!70!black,anchor=west] at (2.1,6.0) {recta con $r$: pendiente $-(1+r)$};
  \node[alert,anchor=west] at (4.6,3.9) {recta con $r'<r$ (más plana)};
  \node[steel,anchor=west] at (0.7,2.75) {inversión en activos reales};
  \draw[{Stealth}-{Stealth},navy] (4.222,-0.95) -- node[below]{inversión} (6,-0.95);
  \draw[{Stealth}-{Stealth},okgreen] (6,-0.95) -- node[below]{VAN} (7.778,-0.95);
\end{tikzpicture}
\end{center}

\paso{2. Quién gana y quién pierde}
\begin{itemize}
  \item \textbf{Prefiere consumo actual} (se endeuda): se mueve hacia el eje horizontal y puede gastar hoy hasta
        $K'>K$. \textbf{Mejora}.
  \item \textbf{Prefiere consumo futuro} (presta): el máximo que puede tener en $t=1$ baja de $M$ a $M'$ porque
        el mercado le remunera menos. \textbf{Empeora}.
\end{itemize}

\begin{solucion}
\textbf{a)}: aumenta la inversión que maximiza el VAN y quien prefiere consumo futuro empeora.
\end{solucion}

% ---------------------------------------------------------------------
\problema{2}{Tipos distintos para prestar y para endeudarse}{Test · Tema 1 · ruptura de la separación}

\begin{enunciado}{2}
Indique cuál de las siguientes afirmaciones es correcta:
\begin{opciones}
  \item Si el tipo de interés del endeudamiento es mayor que el tipo de interés del préstamo, la inversión inicial en
        activos reales que maximizaría el VAN sería la misma tanto desde el punto de vista de un accionista que quisiera
        ser prestamista como desde el punto de vista de un accionista que quisiera ser prestatario.
  \item Si el tipo de interés del endeudamiento es mayor que el tipo de interés del préstamo, un accionista que
        quisiera ser prestamista y otro que quisiera ser prestatario no estarían de acuerdo en cuanto al importe de la
        inversión inicial en activos reales necesaria para maximizar el VAN.
  \item Si el tipo de interés del endeudamiento es mayor que el tipo de interés del préstamo, la línea del tipo de
        interés para quien presta será más inclinada que para quien se endeuda.
  \item Las respuestas a) y c) son correctas.
\end{opciones}
\end{enunciado}

La separación de Fisher exige un mercado perfecto con un único tipo. Si $r_e>r_p$ hay \textbf{dos rectas}:
\begin{itemize}
  \item Quien se \textbf{endeuda} usa la recta más inclinada, $-(1+r_e)$: tangencia en $B$, inversión $C_0=BD$ (menor).
        Por eso c) es falsa.
  \item Quien \textbf{presta} usa la recta más plana, $-(1+r_p)$: tangencia en $A$, inversión $C_0'=AD$ (mayor).
\end{itemize}

\begin{center}
\begin{tikzpicture}[font=\sffamily\scriptsize,x=1.1cm,y=0.68cm]
  \draw[-{Stealth[length=2.2mm]},slate] (0,0) -- (9.8,0) node[below left=2pt]{u.m.\ ($t=0$)};
  \draw[-{Stealth[length=2.2mm]},slate] (0,0) -- (0,8.9) node[above right=-1pt]{u.m.\ ($t=1$)};
  \draw[steel,line width=1.1pt,domain=0.6:6,samples=120,smooth] plot (\x,{2*sqrt(6-\x)});
  \draw[gold,line width=1.1pt] (0,8.033) -- (6.694,0);
  \draw[okgreen,line width=1.1pt] (0,5.833) -- (7.778,0);
  \fill[navy] (5.306,1.667) circle (1.6pt);
  \fill[navy] (4.222,2.667) circle (1.6pt);
  \draw[dashed,slate] (5.306,1.667) -- (5.306,0);
  \draw[dashed,slate] (4.222,2.667) -- (4.222,0);
  \foreach \x/\l in {4.222/A,5.306/B,6/D}{\node[below] at (\x,0) {\l};}
  \node[steel,anchor=west] at (0.7,2.75) {inversión en activos reales};
  \node[gold!70!black,anchor=west] at (2.1,6.0) {quien se endeuda: $-(1+r_e)$};
  \node[okgreen!70!black,anchor=west] at (4.6,3.9) {quien presta: $-(1+r_p)$};
  \draw[{Stealth}-{Stealth},gold!70!black] (5.306,-0.95) -- (6,-0.95) node[right]{$C_0$ (prestatario)};
  \draw[{Stealth}-{Stealth},okgreen] (4.222,-2.0) -- (6,-2.0) node[right]{$C_0'$ (prestamista)};
\end{tikzpicture}
\end{center}

\begin{solucion}
\textbf{b)}: con $r_e>r_p$ prestamistas y prestatarios no se ponen de acuerdo en la inversión que maximiza el VAN.
\end{solucion}

% ---------------------------------------------------------------------
\problema{3}{Ejercicio numérico de Fisher}{Problema · Tema 1 · gráfico + riqueza en $t=0$ y $t=1$}

\begin{enunciado}{3}
Un inversor dispone hoy de un presupuesto de 12.000\eur. El tipo de interés del mercado de capitales a un año es
del 4\,\%. Invierte hoy en activos reales 4.000\eur, lo que le permite obtener un VAN de 2.000\eur, cifra que constituye
el máximo VAN posible. Con esta información, y con un horizonte temporal de un año, se pide:
\begin{opciones}
  \item Representación gráfica de la situación descrita.
  \item Cuantía del flujo de caja que el inversor recibiría al cabo de un año, derivado de su inversión en activos
        reales. \textbf{SOL: 6.240\eur}
  \item Cuantía en la que se incrementaría la riqueza del inversor a consecuencia de la inversión en activos reales,
        si tuviera total preferencia por el consumo futuro frente al consumo actual. \textbf{SOL: 2.080\eur}
  \item Cuantía en la que se incrementaría la riqueza del inversor a consecuencia de la inversión en activos reales,
        si tuviera total preferencia por el consumo actual frente al consumo futuro. \textbf{SOL: 2.000\eur}
\end{opciones}
\end{enunciado}

\paso{Datos}
$D=12.000$; $r=4\,\%$; $C_0=4.000$ $\Rightarrow$ le quedan $12.000-4.000=8.000$; VAN $=2.000$ $\Rightarrow$
$VA=C_0+VAN=6.000$.

\paso{b) Flujo dentro de un año}
\[
VAN=-C_0+\frac{C_1}{1+r}\ \Rightarrow\ 2.000=-4.000+\frac{C_1}{1,04}\ \Rightarrow\ C_1=6.000\cdot1,04=\mathbf{6.240\eur}
\]

\paso{c) Preferencia total por consumo futuro (todo a $t=1$)}
\begin{itemize}
  \item Con inversión real: cobra 6.240 y presta los 8.000 sobrantes: $6.240+8.000\cdot1,04=\mathbf{14.560}$ ($M$).
  \item Solo con el mercado: presta los 12.000: $12.000\cdot1,04=\mathbf{12.480}$ ($L$).
  \item $\Delta$riqueza$_{t=1}=14.560-12.480=\mathbf{2.080\eur}=VAN\cdot(1+r)$.
\end{itemize}

\paso{d) Preferencia total por consumo actual (todo a $t=0$)}
\begin{itemize}
  \item Con inversión real: conserva 8.000 y se endeuda hoy contra los 6.240 futuros ($6.240/1,04=6.000$):
        $8.000+6.000=\mathbf{14.000}$ ($K$).
  \item Sin invertir: 12.000. \ $\Delta$riqueza$_{t=0}=14.000-12.000=\mathbf{2.000\eur}=VAN$.
\end{itemize}

\paso{a) Representación gráfica}
\begin{center}
\begin{tikzpicture}[font=\sffamily\scriptsize,x=0.62cm,y=0.42cm]
  \draw[-{Stealth[length=2.2mm]},slate] (0,0) -- (19.4,0) node[below left=2pt]{u.m.\ ($t=0$, miles)};
  \draw[-{Stealth[length=2.2mm]},slate] (0,0) -- (0,16.3) node[above right=-1pt]{u.m.\ ($t=1$, miles)};
  % recta del mercado tras invertir: de (0;14,56) a (14;0)
  \draw[gold,line width=1.1pt] (0,14.56) -- (14,0);
  % recta del mercado sin invertir: de (0;12,48) a (12;0)
  \draw[slate,dashed,line width=0.8pt] (0,12.48) -- (12,0);
  % curva de inversión y = 2,4763 (12-x)^(2/3), tangente en G(8; 6,24)
  \draw[steel,line width=1.1pt,domain=1.2:12,samples=140,smooth] plot (\x,{2.4763*(12-\x)^(2/3)});
  \fill[navy] (8,6.24) circle (2pt) node[above right=1pt]{G};
  \draw[dashed,slate] (8,6.24) -- (8,0);
  \draw[dashed,slate] (8,6.24) -- (0,6.24);
  \node[left] at (0,14.56) {14.560 (M)};
  \node[left] at (0,12.48) {12.480 (L)};
  \node[left] at (0,6.24) {6.240};
  \node[below] at (8,0) {8.000};
  \node[below] at (12,0) {12.000};
  \node[below] at (14.3,0) {14.000};
  \node[above right=-1pt] at (12,0) {D};
  \node[above right=-1pt] at (14,0) {K};
  \draw[decorate,decoration={brace,mirror,amplitude=4pt},navy] (8,-1.4) -- node[below=4pt]{$C_0=4.000$} (12,-1.4);
  \draw[decorate,decoration={brace,mirror,amplitude=4pt},okgreen] (12,-1.4) -- node[below=4pt]{VAN $=2.000$} (14,-1.4);
  \draw[decorate,decoration={brace,mirror,amplitude=4pt},slate] (8,-3.6) -- node[below=4pt]{$VA=6.000$} (14,-3.6);
  \draw[decorate,decoration={brace,amplitude=4pt},okgreen] (-3.3,12.48) -- node[left=4pt,align=right]{$\Delta$riq.\ $t{=}1$\\$=2.080$} (-3.3,14.56);
  \draw[steel,-{Stealth[length=1.8mm]}] (3.2,4.3) -- (3.55,9.9);
  \node[steel,below] at (3.2,4.3) {inversión en activos reales};
  \node[gold!70!black,anchor=west] at (10.2,5.2) {mercado, $-(1{,}04)$};
\end{tikzpicture}
\end{center}

\begin{clave}
El VAN es el aumento de riqueza medido en $t=0$; medido en $t=1$ es $VAN\cdot(1+r)$. Los dos inversores mejoran,
sea cual sea su preferencia de consumo: eso es la separación de Fisher.
\end{clave}

% =====================================================================
\bloque{Tema 2}{Valoración de acciones y VAOC}{Problemas 4 a 11 · Gordon, $g=b\cdot ROE$, crecimiento por etapas y oportunidades de crecimiento}
% =====================================================================

\begin{clave}[Herramientas del bloque]
$P_0=\dfrac{DIV_1}{r-g}$ \quad $g=b\cdot ROE$ \quad $DIV=(1-b)\,BPA$ \quad
$P_0=\dfrac{BPA_1}{r}+\text{VAOC}$ \quad $VAN_t=-I_t+\dfrac{ROE\cdot I_t}{r}$.\\
Cuando se agotan las oportunidades de crecer, $DIV=BPA$ constante y $P_T=BPA_{T+1}/r$.
\end{clave}

\problema{4}{Crecimiento con ROE igual a \texorpdfstring{$r$}{r}}{Test · Gordon y VAOC}

\begin{enunciado}{4}
Una empresa que el próximo año espera obtener un BPA de 20 euros/acción prevé que sus dividendos pueden crecer
indefinidamente al 4\,\% anual. Si el dividendo del próximo año es de 10 euros/acción y la tasa de capitalización del
mercado es del 8\,\%, podemos afirmar que:
\begin{opciones}
  \item Esta empresa tiene oportunidades de crecimiento que el mercado valora en 250 euros/acción, debido a que
        reinvierte parte de sus beneficios y obtiene de esas reinversiones una tasa de rentabilidad conveniente.
  \item Todo lo que el mercado valora en esta empresa son sus oportunidades de crecimiento (250 euros/acción).
  \item La empresa tiene oportunidades de crecimiento, ya que reinvierte parte de sus beneficios, pero no es posible
        cuantificarlas, al no conocerse la rentabilidad que está obteniendo en sus reinversiones.
  \item El VAOC de esta empresa es nulo pues, pese a reinvertir parte de sus beneficios, la rentabilidad que obtiene en
        sus reinversiones es exactamente igual a la tasa de capitalización del mercado.
\end{opciones}
\end{enunciado}

\[
P_0=\frac{10}{0,08-0,04}=250 \qquad d=\frac{10}{20}=0,5\Rightarrow b=0,5 \qquad
ROE=\frac{g}{b}=\frac{0,04}{0,5}=8\,\%=r
\]
\[
\text{VAOC}=P_0-\frac{BPA_1}{r}=250-\frac{20}{0,08}=250-250=\mathbf{0}
\]

\begin{solucion}
\textbf{d)}: crece, pero sus reinversiones rinden justo $r$; crecer sin superar el coste del capital no crea valor.
\end{solucion}

% ---------------------------------------------------------------------
\problema{5}{Crecimiento hasta el año 6}{Test · descuento de dividendos en dos etapas}

\begin{enunciado}{5}
Una empresa que el año próximo prevé obtener un BPA de 25 euros/acción tiene una tasa de reparto de dividendos
prevista del 40\,\%. La empresa estima que sus dividendos podrán crecer hasta el año 6 a razón de un 6\,\% anual. Del
año 6 en adelante, se prevé que sus oportunidades de crecer se habrán agotado, por lo que se repartirán
indefinidamente todos sus beneficios como dividendos. Si la tasa de capitalización del mercado es del 9\,\%, el precio
de una acción de esta empresa (utilizando 2 decimales para los cálculos) es de:
\begin{opciones}
  \item 265,09 euros/acción.
  \item 285,04 euros/acción.
  \item 415,19 euros/acción.
  \item 333,33 euros/acción.
\end{opciones}
\end{enunciado}

El BPA crece al 6\,\% hasta el año 6 ($BPA_6=25\cdot1,06^5=33,46$). Desde el año 6 se reparte todo:
$DIV_6=DIV_7=\dots=33,46$.

\begin{center}\small
\begin{tabularx}{0.86\textwidth}{@{}L RRRRRR@{}}
\headrow \thead{€/acción} & \thead{Año 1} & \thead{Año 2} & \thead{Año 3} & \thead{Año 4} & \thead{Año 5} & \thead{Año 6 y sig.}\\
$BPA_t$ & 25,00 & 26,50 & 28,09 & 29,78 & 31,56 & 33,46\\
\rowcolor{mist} Reparto & 40\,\% & 40\,\% & 40\,\% & 40\,\% & 40\,\% & 100\,\%\\
$DIV_t$ & 10,00 & 10,60 & 11,24 & 11,91 & 12,62 & 33,46\\
\end{tabularx}
\end{center}

\[
P_5=\frac{DIV_6}{r}=\frac{33,46}{0,09}=371,78
\qquad
P_0=\frac{10}{1,09}+\frac{10,60}{1,09^2}+\frac{11,24}{1,09^3}+\frac{11,91}{1,09^4}+\frac{12,62+371,78}{1,09^5}=\mathbf{285,04}
\]

\begin{solucion}
\textbf{b) 285,04 €/acción}.
\end{solucion}
\begin{trampa}[Redondeo]
Con todos los decimales sale 285,01. El enunciado pide trabajar con 2 decimales, y así se obtiene exactamente 285,04.
\end{trampa}

% ---------------------------------------------------------------------
\problema{6}{Obtener \texorpdfstring{$g$ y $r$}{g y r} a partir del precio}{Test · Gordon despejando $r$}

\begin{enunciado}{6}
Las acciones de la sociedad X cotizan actualmente a 73 euros/acción. El BPA esperado para el próximo año es de 2,80
euros/acción, del que se espera repartir un dividendo de 1,68 euros/acción. Si la rentabilidad que esta empresa puede
esperar obtener en sus reinversiones es del 20\,\%, es posible afirmar que:
\begin{opciones}
  \item La tasa esperada de crecimiento del dividendo será del 12\,\%, y la tasa de capitalización del mercado
        ascenderá aproximadamente al 14,3\,\%.
  \item La tasa esperada de crecimiento del dividendo será del 8\,\%, y la tasa de capitalización del mercado
        ascenderá aproximadamente al 10,3\,\%.
  \item La tasa esperada de crecimiento del dividendo será del 10,3\,\%, y la tasa de capitalización del mercado
        ascenderá aproximadamente al 8\,\%.
  \item Las respuestas a), b) y c) son falsas.
\end{opciones}
\end{enunciado}

\[
d=\frac{1,68}{2,80}=0,60\Rightarrow b=0,40 \qquad g=0,40\cdot0,20=\mathbf{8\,\%} \qquad
r=\frac{1,68}{73}+0,08=0,0230+0,08=\mathbf{10,3\,\%}
\]

\begin{solucion}
\textbf{b)}: $g=8\,\%$ y $r\approx10,3\,\%$.
\end{solucion}
\begin{trampa}[Distractor típico]
La opción a) usa $g=0,60\cdot0,20=12\,\%$: aplica la tasa de \emph{reparto} en lugar de la de \emph{retención}.
\end{trampa}

% ---------------------------------------------------------------------
\problema{7}{De qué depende el crecimiento}{Test · teoría de Gordon}

\begin{enunciado}{7}
Indique cuál de las siguientes afirmaciones es correcta:
\begin{opciones}
  \item La fórmula del flujo de tesorería descontado con crecimiento constante es apropiada para calcular el precio de
        las acciones de una empresa que viene experimentando en los últimos años una tasa muy elevada de crecimiento de
        los dividendos.
  \item El crecimiento que una empresa pueda experimentar en el futuro depende exclusivamente de la proporción de sus
        beneficios anuales que la misma destine a reinversión.
  \item El crecimiento que una empresa pueda experimentar en el futuro depende exclusivamente de la rentabilidad que
        pueda obtener de las oportunidades de inversión futuras que la misma encuentre.
  \item Las respuestas a), b) y c) son falsas.
\end{opciones}
\end{enunciado}

\begin{itemize}
  \item \textbf{a) es falsa}: un crecimiento muy alto suele ser transitorio y puede incluso superar a $r$. Aplicar
        $DIV_1/(r-g)$ con ese $g$ sobrevalora la acción (o no tiene sentido si $g\ge r$): hay que usar etapas.
  \item \textbf{b) y c) son falsas}: $g=b\cdot ROE$ depende de \textbf{las dos} cosas a la vez, no de una sola.
\end{itemize}

\begin{solucion}
\textbf{d)}: las tres afirmaciones son falsas.
\end{solucion}

% ---------------------------------------------------------------------
\problema{8}{MILLION, S.A.: una oportunidad puntual en \texorpdfstring{$t=2$}{t=2}}{Problema · VAOC de una única reinversión}

\begin{enunciado}{8}
La empresa MILLION, S.A. espera obtener el próximo año un BPA de 20\eur/acción. Su director financiero estima que dentro
de dos años ($t=2$) surgirá una oportunidad de inversión interesante que piensa que podría reportar una rentabilidad
anual a perpetuidad sobre el capital invertido del 15\,\%. La tasa de capitalización del mercado es del 12\,\%.
Con esta información, indique:
\begin{opciones}
  \item Si se decidiera reinvertir el 80\,\% de los beneficios esperados por la empresa en el año 2, ¿cuál sería el
        valor actual de las oportunidades de crecimiento en cada acción de la empresa? \textbf{SOL: VAOC = 3,19\eur}
  \item ¿Cuál sería el precio actual de la acción? \textbf{SOL: $P_0$ = 169,86\eur}
  \item ¿Cuál es la rentabilidad anual mínima que usted consideraría aceptable para la reinversión que la empresa se
        plantea realizar en $t=2$? Justifique su respuesta.
\end{opciones}
\end{enunciado}

\paso{a) VAOC}
Sin reinversión en $t=1$, $BPA_2=BPA_1=20$. En $t=2$ se reinvierte $I_2=0,80\cdot20=16$, que rinde $16\cdot0,15=2,40$
al año a perpetuidad desde $t=3$.
\[
VAN_2=-16+\frac{2,40}{0,12}=-16+20=4 \qquad
\text{VAOC}=\frac{4}{1,12^2}=\mathbf{3,19\eur}
\]

\paso{b) Precio}
\[
P_0=\frac{BPA_1}{r}+\text{VAOC}=\frac{20}{0,12}+3,19=166,67+3,19=\mathbf{169,86\eur}
\]
Comprobación por dividendos: $DIV_1=20$, $DIV_2=4$ y desde $t=3$, $DIV=22,40$:
\(
P_0=\frac{20}{1,12}+\frac{4+22,40/0,12}{1,12^2}=17,86+152,00=169,86.
\)

\paso{c) Rentabilidad mínima}
El coste de oportunidad, \textbf{12\,\%}. Con $ROE=12\,\%$: $VAN_2=-16+16\cdot0,12/0,12=0$ (indiferente). Por encima,
la reinversión crea valor; por debajo, conviene repartir el dividendo.

\begin{solucion}
VAOC $=3,19$\eur; \ $P_0=169,86$\eur; \ rentabilidad mínima exigible $=r=12\,\%$.
\end{solucion}

% ---------------------------------------------------------------------
\problema{9}{Empresa Z: crecimiento en los años 1, 2 y 3}{Problema · multietapa por DDM y por VAOC}

\begin{enunciado}{9}
Disponemos de la siguiente información de la empresa Z:
\begin{itemize}
  \item $BPA_1=80$ euros/acción.
  \item Tasa de reparto de dividendos: 80\,\%.
  \item Se espera que los beneficios crezcan hasta el año 4, gracias a la obtención de una rentabilidad del 20\,\% en sus
        reinversiones de los años 1, 2 y 3.
  \item Desde ese año en adelante, se consideran agotadas las oportunidades de crecimiento, de forma que se repartirán
        todos los beneficios como dividendos.
  \item La tasa de capitalización del mercado es del 12\,\%.
\end{itemize}
SE PIDE:
\begin{opciones}
  \item El precio de la acción de la empresa Z en el año 4. \textbf{SOL: $P_4$ = 749,92\eur}
  \item El precio actual de la acción de la empresa Z. \textbf{SOL: $P_0$ = 693,25\eur}
\end{opciones}
\end{enunciado}

\paso{Crecimiento y cuadro}
$b=1-0,80=0,20$; $g=0,20\cdot0,20=4\,\%$ hasta el año 4. Desde el año 4 se reparte el 100\,\%.

\begin{center}\small
\begin{tabularx}{0.8\textwidth}{@{}L RRRR@{}}
\headrow \thead{€/acción} & \thead{Año 1} & \thead{Año 2} & \thead{Año 3} & \thead{Año 4 y sig.}\\
$BPA_t$ & 80,00 & 83,20 & 86,53 & 89,99\\
\rowcolor{mist} Reparto & 80\,\% & 80\,\% & 80\,\% & 100\,\%\\
$DIV_t$ & 64,00 & 66,56 & 69,22 & 89,99\\
\rowcolor{mist} Reinversión $I_t$ & 16,00 & 16,64 & 17,31 & 0\\
\end{tabularx}
\end{center}

\paso{a) Precio en el año 4}
\[
P_4=\frac{DIV_5}{r}=\frac{89,99}{0,12}=\mathbf{749,92\eur}
\]

\paso{b) Precio hoy por descuento de dividendos}
\[
P_0=\frac{64}{1,12}+\frac{66,56}{1,12^2}+\frac{69,22}{1,12^3}+\frac{89,99+749,92}{1,12^4}
=57,14+53,06+49,27+533,78=\mathbf{693,25\eur}
\]

\paso{Comprobación por VAOC}
$VAN_1=-16+\frac{3,20}{0,12}=10,67$; \ $VAN_2=-16,64+\frac{3,328}{0,12}=11,09$; \
$VAN_3=-17,31+\frac{3,461}{0,12}=11,54$
\[
\text{VAOC}=\frac{10,67}{1,12}+\frac{11,09}{1,12^2}+\frac{11,54}{1,12^3}=26,58
\qquad P_0=\frac{80}{0,12}+26,58=666,67+26,58=693,25
\]

\begin{solucion}
$P_4=749,92$\eur; \ $P_0=693,25$\eur.
\end{solucion}

% ---------------------------------------------------------------------
\problema{10}{Green Rock Corp.: la competencia reduce la ROE}{Problema · dos etapas por DDM y por VAOC}

\begin{enunciado}{10}
La empresa Green Rock Corp., tiene una rentabilidad sobre el capital invertido del 10\,\% a perpetuidad y una tasa de
crecimiento constante de sus reinversiones del 5\,\%. El valor contable por acción es actualmente de 91\eur. Suponga que
la ROE y la tasa de reparto de dividendos se mantiene constante durante los próximos 2 años. Después de esa fecha, la
competencia hará que la ROE descienda al 8\,\% a perpetuidad y la tasa de crecimiento constante baje al 2,4\,\%. La tasa
de capitalización del mercado se estima en el 8\,\%. Con esta información, determine:
\begin{enumerate}
  \item El precio actual de la acción de la empresa Green Rock Corp., empleando el método del descuento de dividendos.
  \item El precio actual de la acción de la empresa, detallando el cálculo previo del Valor Actual de la Oportunidades
        de Crecimiento (VAOC), explicando qué parte del precio se debe a las oportunidades de crecimiento.
\end{enumerate}
\end{enunciado}

\paso{Parámetros de cada etapa}
\begin{itemize}
  \item Años 1--2: $ROE=10\,\%$, $g=5\,\%$ $\Rightarrow b=0,05/0,10=50\,\%$, reparto 50\,\%.
  \item Año 3 en adelante: $ROE=8\,\%=r$, $g=2,4\,\%$ $\Rightarrow b=0,024/0,08=30\,\%$, reparto 70\,\%.
\end{itemize}

\begin{center}\small
\begin{tabularx}{0.82\textwidth}{@{}c RcRR@{}}
\headrow \thead{Año} & \thead{$BPA_t$} & \thead{Reparto} & \thead{$DIV_t$} & \thead{Reinversión $I_t$}\\
1 & $91\cdot0,10=9,1000$ & 50\,\% & 4,5500 & 4,5500 (al 10\,\%)\\
\rowcolor{mist} 2 & $9,10\cdot1,05=9,5550$ & 50\,\% & 4,7775 & 4,7775 (al 10\,\%)\\
3 & $9,555\cdot1,05=10,0328$ & 70\,\% & 7,0229 & 3,0098 (al 8\,\%)\\
\rowcolor{mist} 4 & $10,0328\cdot1,024=10,2735$ & 70\,\% & 7,1915 & 3,0821 (al 8\,\%)\\
\end{tabularx}
\end{center}

\paso{1. Descuento de dividendos}
Desde el año 3 el dividendo crece al 2,4\,\% indefinidamente, así que
$P_2=\dfrac{DIV_3}{r-g}=\dfrac{7,0229}{0,08-0,024}=125,41$:
\[
P_0=\frac{4,55}{1,08}+\frac{4,7775+125,41}{1,08^2}=4,21+111,62=\mathbf{115,83\eur}
\]

\paso{2. Por VAOC}
\[
VAN_1=-4,55+\frac{4,55\cdot0,10}{0,08}=1,1375 \qquad VAN_2=-4,7775+\frac{4,7775\cdot0,10}{0,08}=1,1944
\]
Desde el año 3 la ROE es igual a $r$, así que esas reinversiones tienen $VAN=0$.
\[
\text{VAOC}=\frac{1,1375}{1,08}+\frac{1,1944}{1,08^2}=1,05+1,02=2,08
\qquad
P_0=\frac{9,10}{0,08}+2,08=113,75+2,08=\mathbf{115,83\eur}
\]

\begin{solucion}
$P_0=115,83$\eur por los dos métodos. Las oportunidades de crecimiento explican solo \textbf{2,08\eur (1,8\,\%)};
el 98,2\,\% (113,75\eur) es el valor de los activos actuales sin crecimiento.
\end{solucion}

% ---------------------------------------------------------------------
\problema{11}{North Atlantic, S.A.: el proyecto de la avena}{Problema · VAOC de un proyecto y ratio $BPA_1/P_0$}

\begin{enunciado}{11}
La empresa North Atlantic S.A. se dedica al cultivo y venta de trigo, lo que le viene reportando unos beneficios netos
anuales de 4.500.000 euros, importe que se prevé estable para los próximos años. Actualmente se está planteando
introducirse en el cultivo y distribución de avena, lo que supondrá una inversión de 1.000.000 de euros en el año 1,
previéndose la venta de 10.000 Tm/año de avena desde el año 2 en adelante. Los costes de producción y distribución
ascienden a 100\eur/Tm, importe que se mantendrá constante a futuro y el precio previsto para la avena en los próximos años
se estima en: año 1: 140\eur/Tm; año 2 y siguientes: 150\eur/Tm.

North Atlantic S.A. se financia sólo con fondos propios: tiene 1.000.000 de acciones en circulación y el coste de
capitalización de las acciones es del 10\,\%. Todos los beneficios de la empresa se reparten como dividendos excepto en
el año 1, en que se reinvertirá la cantidad necesaria para llevar a cabo el nuevo proyecto de la avena. Suponga
asimismo que la empresa mantendrá su actividad a perpetuidad y que estamos en un contexto sin impuestos.
\begin{opciones}
  \item Determine el valor actual de las oportunidades de crecimiento de esta empresa.
  \item Calcule el precio actual de las acciones de esta empresa.
  \item Calcule el ratio $BPA_1/P_0$ y explique por qué no coincide con el coste de capital del 10\,\%.
\end{opciones}
\end{enunciado}

\paso{Datos por acción (1.000.000 de acciones)}
$BPA_1=4,50$; inversión en $t=1$: 1,00 por acción; flujo de la avena desde $t=2$:
$10.000\cdot(150-100)=500.000$\eur al año, es decir, 0,50 por acción. El precio del año 1 (140) no se usa porque ese
año todavía no se vende avena.

\begin{center}
\flujos[2.0]{4}{1/n/{$-1{,}00$ (inversión)},2/p/{$+0{,}50$},3/p/{$+0{,}50$},4/p/{$+0{,}50$ \dots}}
\end{center}

\paso{a) VAOC}
\[
VAN_1=-1+\frac{0,50}{0,10}=4,00 \qquad \text{VAOC}=\frac{4,00}{1,10}=\mathbf{3,64}\text{\,€/acción}
\quad(3.636.363,64\eur\text{ en total})
\]

\paso{b) Precio}
\[
P_0=\frac{4,50}{0,10}+3,64=45,00+3,64=\mathbf{48,64\eur}
\]
Comprobación por dividendos: $DIV_1=4,50-1,00=3,50$; desde $t=2$, $DIV=5,00$:
$P_0=\frac{3,50+5,00/0,10}{1,10}=\frac{53,50}{1,10}=48,64$.

\paso{c) Ratio beneficio/precio}
\[
\frac{BPA_1}{P_0}=\frac{4,50}{48,64}=\mathbf{9,25\,\%}<10\,\% \qquad
\text{porque}\quad \frac{BPA_1}{P_0}=r\left(1-\frac{\text{VAOC}}{P_0}\right)
\]

\begin{solucion}
VAOC $=3,64$\eur; \ $P_0=48,64$\eur; \ $BPA_1/P_0=9,25\,\%$. Solo coincide con $r$ si VAOC $=0$; aquí la avena
rinde un 50\,\% $>10\,\%$ y el precio incorpora ese crecimiento.
\end{solucion}

% =====================================================================
\bloque{Tema 3}{Criterios de valoración}{Problema 12 · VAN frente a TIR}
% =====================================================================

\problema{12}{VAN frente a TIR}{Test · proyectos mutuamente excluyentes}

\begin{enunciado}{12}
Indique cuál de las siguientes afirmaciones sobre algunos de los criterios de selección de inversiones es correcta:
\begin{opciones}
  \item El criterio TIR resulta adecuado para elegir entre dos proyectos de inversión que tienen distinto desembolso inicial.
  \item El criterio VAN es adecuado para elegir entre dos proyectos que tienen el mismo horizonte temporal.
  \item El criterio TIR es adecuado para elegir entre dos proyectos de inversión en cualquier caso.
  \item En cualquier circunstancia los criterios VAN y TIR llevarán a la misma decisión.
\end{opciones}
\end{enunciado}

\begin{itemize}
  \item \textbf{a), c) y d) son falsas}: la TIR es una medida relativa. Ignora la escala del desembolso y el reparto
        temporal de los flujos, no distingue entre prestar y endeudarse y puede dar varias tasas o ninguna. VAN y TIR
        pueden ordenar los proyectos al revés cuando $r$ está por debajo de la intersección de Fisher (salvo que se
        calcule la TIR del proyecto incremental).
  \item \textbf{b) es correcta}: el VAN mide en euros de hoy la riqueza creada y sirve para elegir entre proyectos
        excluyentes con el mismo horizonte. Con horizontes distintos y equipos que se reponen, se usa el CAE (P23).
\end{itemize}

\begin{solucion}
\textbf{b)}.
\end{solucion}

% =====================================================================
\bloque{Tema 4}{Flujos de caja y decisiones de inversión}{Problemas 13 a 25 · $C_0\di$, $Q_i\di$, $\Delta FM$, $VR_n\di$, análisis incremental, CAE y demanda estacional}
% =====================================================================

\problema{13}{Conservas de pescado: la nave tiene coste de oportunidad}{Test · VAN incremental con FM no recuperado}

\begin{enunciado}{13}
Una fábrica de conservas de pescado pretende acometer un proyecto de ampliación de su planta envasadora en una nave de
su propiedad, que actualmente se encuentra sin utilización específica. La inversión necesaria en equipos es de
900.000\eur, que se amortizarán linealmente en 10 años. Este proyecto provocará un incremento inicial del fondo de
maniobra de 12.000\eur. Los ingresos anuales se estiman, para los 10 años de vida útil de los equipos, en 1.200.000\eur. Los
costes operativos anuales serán previsiblemente de 535.000\eur. Se estima que en cada año el fondo de maniobra ascenderá
al 2\,\% de las ventas anuales previstas. El valor residual de los equipos se puede considerar despreciable al final de
su vida útil. El fondo de maniobra no se recupera al final de la vida del proyecto. De no acometerse el proyecto, la
nave se alquilaría por 360.000\eur al año durante 10 años.

Sabiendo que el tipo del impuesto de sociedades es del 35\,\%, que toda la cuota de amortización contable es deducible a
efectos fiscales, y que la tasa de actualización apropiada para reflejar el riesgo del proyecto es del 6,5\,\%, indique
cuál de las siguientes afirmaciones es correcta:
\begin{opciones}
  \item Se aconseja no acometer el proyecto de ampliación, pues el VAN incremental (de acometer el proyecto frente a no
        acometerlo) arroja una cifra de $-732.841,04$\eur.
  \item Se aconseja acometer el proyecto, pues el VAN incremental (de acometer el proyecto frente a no acometerlo)
        arroja una cifra de 728.366,14\eur.
  \item Se aconseja acometer el proyecto, pues el VAN incremental (de acometer el proyecto frente a no acometerlo)
        asciende a 732.841,04\eur.
  \item Las respuestas a), b) y c) son falsas.
\end{opciones}
\end{enunciado}

\paso{1. Desembolso inicial}
$\Delta C_0\di=-900.000-12.000=-912.000$ (la nave es propia: no se paga, pero se renuncia a su alquiler).

\paso{2. Fondo de maniobra}
$FM_0=12.000$; \ $FM_{1\ldots10}=2\,\%\cdot1.200.000=24.000$ $\Rightarrow$ $\Delta FM_1=+12.000$ y
$\Delta FM_{2\ldots10}=0$.

\paso{3. Flujos de las operaciones}
El alquiler perdido es un \textbf{coste de oportunidad}: se resta como un coste más. $CAF=900.000/10=90.000$.
\[
\Delta Q_i\ai=1.200.000-535.000-360.000=305.000
\qquad
\Delta Q_i\di\text{ renta}=305.000\cdot0,65+90.000\cdot0,35=\mathbf{229.750}
\]
Tesorería: año 1: $229.750-12.000=217.750$; años 2--10: $229.750$.

\paso{4. Valor residual}
Los 24.000 de FM no se recuperan: se pierde su valor contable y esa pérdida es deducible,
$VR_{10}\di=0+24.000\cdot0,35=8.400$.

\begin{center}
\flujos[1.25]{10}{0/n/{$-912.000$},1/p/{217.750},2/p/{229.750},5/p/{$\cdots$},9/p/{229.750},10/p/{229.750\\$+$8.400}}
\end{center}

\paso{5. VAN incremental ($r=6,5\,\%$, \ $\an{10}{0{,}065}=7,188830$)}
\begin{align*}
\Delta VAN&=-912.000-\frac{12.000}{1,065}+229.750\cdot7,188830+\frac{8.400}{1,065^{10}}\\
&=-912.000-11.267,61+1.651.633,75+4.474,90=\mathbf{732.841,04}
\end{align*}

\begin{solucion}
\textbf{c)} $\Delta VAN=732.841,04$\eur $>0$: conviene ampliar.
\end{solucion}
\begin{trampa}[Por qué existe la opción b)]
Si se olvida el ahorro fiscal por la pérdida del FM no recuperado ($8.400$ en $t=10$), sale 728.366,14 (opción b).
La diferencia es justo $8.400/1,065^{10}=4.474,90$.
\end{trampa}

% ---------------------------------------------------------------------
\problema{14}{Reglas del desembolso inicial}{Test · costes hundidos, renovación y FM}

\begin{enunciado}{14}
Indique cuál de las siguientes afirmaciones sobre la valoración de inversiones con el criterio del VAN es
\textbf{falsa}:
\begin{opciones}
  \item Al calcular el desembolso inicial de un proyecto de inversión no habrá que incorporar el importe de los
        desembolsos pasados e irreversibles, aunque éstos se efectuaran para posibilitar la realización del proyecto.
  \item En una inversión de renovación, el importe percibido por la venta del equipo que se sustituye aumentará la
        cuantía del desembolso inicial incremental.
  \item En una inversión de renovación, el importe percibido por la venta del equipo que se sustituye reducirá la
        cuantía del desembolso inicial incremental.
  \item Al calcular el desembolso inicial habrá que considerar las variaciones en el fondo de maniobra de la empresa
        causadas por el proyecto.
\end{opciones}
\end{enunciado}

La venta del equipo viejo es un \textbf{cobro} en $t=0$: \textbf{reduce} el desembolso incremental,
$\Delta C_0\di=C_0\text{ nuevo}-VR_0\di\text{ viejo}$. Por eso c) es verdadera y b) falsa. a) (costes hundidos fuera)
y d) (el FM inicial sí cuenta) son verdaderas.

\begin{solucion}
La afirmación falsa es la \textbf{b)}.
\end{solucion}

% ---------------------------------------------------------------------
\problema{15}{Renovación de una planta}{Test · qué entra en un análisis incremental}

\begin{enunciado}{15}
Determinada empresa actualmente utiliza una planta de producción que le genera unos flujos de caja de 1 millón de\eur
anuales durante 5 años. Si la planta se vendiese en el momento actual obtendría por ella 6 millones de\eur y si lo
hiciera dentro de 5 años recibiría 2 millones de\eur. La empresa está planteando renovar la planta por otro equipo que le
costaría hoy 8 millones de\eur y que le producirían 2,5 millones de\eur durante los 5 años, con valor residual nulo. Si se
plantea estudiar la viabilidad financiera de la renovación, mediante el criterio del VAN incremental, se puede afirmar
que:
\begin{opciones}
  \item El coste de la renovación no debe considerar la venta de la planta actual.
  \item El coste de la renovación sí debe considerar la venta de la planta actual, pero el flujo de caja de la
        renovación no debe considerar los flujos de caja proporcionados por la planta antigua.
  \item En este tipo de proyectos, el coste de la renovación no debe considerar la venta de la planta actual, pero el
        flujo de caja de la renovación si debe considerar los flujos de caja proporcionados por la planta antigua.
  \item Ninguna de las opciones anteriores (a, b y c) son correctas.
\end{opciones}
\end{enunciado}

En el análisis incremental (renovar $-$ seguir) \textbf{todo} lo de la planta vieja entra con signo contrario:
\begin{center}\small
\begin{tabularx}{0.78\textwidth}{@{}L QQQ@{}}
\headrow \thead{Millones de €} & \thead{$t=0$} & \thead{$t=1\ldots5$} & \thead{$t=5$ (VR)}\\
Renovar (planta nueva) & $-8$ & $+2,5$ & 0\\
\rowcolor{mist} Seguir con la vieja (se pierde) & $+6$ & $-1,0$ & $-2$\\
\totrow \textbf{Incremental} & $\mathbf{-2}$ & $\mathbf{+1,5}$ & $\mathbf{-2}$\\
\end{tabularx}
\end{center}

\begin{solucion}
\textbf{d)}: hay que considerar la venta de la planta actual \textbf{y} los flujos que dejaría de generar (y su valor
residual), así que a), b) y c) son falsas.
\end{solucion}

% ---------------------------------------------------------------------
\problema{16}{Bombillas: punto muerto financiero}{Problema · cantidad mínima con VAN $=0$}

\begin{enunciado}{16}
Una empresa pretende adquirir un nuevo equipo para fabricar bombillas de bajo consumo, cuyo coste inicial asciende a
800.000\eur, y que se amortizará linealmente en 10 años. Se estima que el precio de venta será de 12\eur/bombilla; los costes
variables se prevé que ascenderán a 6,50\eur/bombilla, y los costes fijos se estiman en 150.000\eur/año (esta cifra incluye
la amortización técnica prevista del equipo). La duración del proyecto será de 10 años, al final de los cuales el valor
residual del equipo se puede despreciar. El tipo vigente para el impuesto de sociedades es del 35\,\%, y la tasa de
actualización que refleja adecuadamente el riesgo de la inversión es del 8\,\%. Indique la cantidad mínima de bombillas
que la empresa debería vender para que resultara aconsejable acometer el proyecto. \textbf{SOL: 38.243 bombillas}
\end{enunciado}

\paso{1. Separar la amortización de los costes fijos}
$CAF=800.000/10=80.000$ $\Rightarrow$ costes fijos que se pagan $=150.000-80.000=70.000$. Margen unitario
$=12-6,50=5,50$.

\paso{2. Flujo anual en función de las unidades $X$}
\[
Q\di=(5,50X-70.000)\cdot0,65+80.000\cdot0,35=3,575X-17.500
\]

\paso{3. VAN $=0$ ($\an{10}{0{,}08}=6,710081$)}
\[
-800.000+(3,575X-17.500)\cdot6,710081=0\ \Rightarrow\ Q\di=119.223,59\ \Rightarrow\ X=\frac{136.723,59}{3,575}=38.244,36
\]

\begin{center}
\begin{tikzpicture}
\begin{axis}[width=0.78\textwidth,height=5.2cm,
  xlabel={\sffamily\scriptsize bombillas vendidas al año},ylabel={\sffamily\scriptsize VAN (miles de €)},
  xmin=20000,xmax=56000,ymin=-480,ymax=480,
  axis lines=left,tick label style={font=\sffamily\scriptsize},
  scaled x ticks=false,xtick={20000,30000,38244,46000,56000},
  xticklabels={20.000,30.000,\textbf{38.244},46.000,56.000},
  ytick={-400,-200,0,200,400},
  grid=major,grid style={rule!60},
  every axis plot/.append style={line width=1.2pt}]
  \addplot[steel,domain=20000:56000,samples=2] {(-800000+(3.575*x-17500)*6.710081)/1000};
  \addplot[alert,dashed,line width=0.8pt] coordinates {(20000,0) (56000,0)};
  \addplot[only marks,mark=*,mark size=2.4pt,navy] coordinates {(38244.36,0)};
  \node[font=\sffamily\scriptsize,okgreen!70!black,anchor=west] at (axis cs:44000,250) {VAN $>0$: conviene};
  \node[font=\sffamily\scriptsize,alert,anchor=west] at (axis cs:21000,-300) {VAN $<0$: no conviene};
\end{axis}
\end{tikzpicture}
\end{center}

\begin{solucion}
Hacen falta unas \textbf{38.243} bombillas al año (solución oficial).
\end{solucion}
\begin{trampa}[Redondeo de la solución oficial]
Con el factor exacto sale $X=38.244,36$, así que el entero que da VAN $\ge0$ es \textbf{38.245}. La cifra oficial
(38.243) sale de redondear el factor de actualización. En el examen escribe tu cálculo completo: la diferencia es
solo de redondeo.
\end{trampa}

% ---------------------------------------------------------------------
\problema{17}{Cuadro de siete periodos}{Problema · beneficios, $C_0\di$, $Q_i\di$ tesorería y $VR\di$}

\begin{enunciado}{17}
Determinada empresa a inicio del año 2023 presenta la información financiera que se recoge en el cuadro adjunto,
relativa a un proyecto de inversión que quiere emprender:

\begingroup\scriptsize\setlength{\tabcolsep}{2.6pt}
\begin{tabularx}{\linewidth}{@{}>{\raggedright\arraybackslash}p{2.75cm} *{8}{R}@{}}
\headrow \thead{Concepto (€)} & \thead{0} & \thead{1} & \thead{2} & \thead{3} & \thead{4} & \thead{5} & \thead{6} & \thead{7}\\
Inversión de capital & 15.000 & & & & & & & \\
Valor residual de venta & & & & & & & & 2.000\\
Cuota de amortización & & 2.250 & 2.250 & 2.250 & 2.250 & 2.250 & 2.250 & \\
Fondo amort.\ acumulada & & 2.250 & 4.500 & 6.750 & 9.000 & 11.250 & 13.500 & \\
Valor contable fin año & & 12.750 & 10.500 & 8.250 & 6.000 & 3.750 & 1.500 & \\
Fondo de maniobra & -- & 1.000 & 500 & 300 & 1.500 & 1.500 & 1.000 & 500\\
Valor contable activo total & & 13.750 & 11.000 & 8.550 & 7.500 & 5.250 & 2.500 & \\
Ingresos por ventas & -- & 1.500 & 5.000 & 30.000 & 50.000 & 35.000 & 30.000 & \\
Coste de las ventas & -- & 837 & 7.729 & 19.552 & 29.345 & 21.492 & 11.830 & \\
Otros costes & 4.000 & 2.200 & 1.200 & 1.300 & 1.500 & 1.600 & 1.700 & \\
\end{tabularx}
\endgroup

\smallskip
La amortización de la inversión inicial es lineal en 6 años, reconociéndose un valor residual contable de 1.500\eur el
sexto año. La tasa impositiva es del 35\,\% anual. El coste de oportunidad del capital que la empresa exige a sus
proyectos del 15\,\%. Suponga que la empresa cuenta con otros proyectos que generan beneficios y que pagan impuestos en
los próximos 10 años y que el fondo de maniobra se recupera por su valor contable. Con estos datos se pide calcular:
\begin{opciones}
  \item Los beneficios antes y después de impuestos.
  \item La inversión inicial después de impuestos, los flujos netos de caja después de impuestos y el valor residual
        después de impuestos.
  \item El VAN del proyecto. \textbf{SOL: VAN = 2.317,57\eur}
\end{opciones}
\end{enunciado}

\paso{a) Beneficios antes y después de impuestos}
$CAF=(15.000-1.500)/6=2.250$. En el periodo 7 la venta por 2.000 de un activo con valor contable 1.500 da una
plusvalía de 500. Las pérdidas generan ahorro fiscal porque la empresa tributa por otros proyectos.

\begingroup\scriptsize\setlength{\tabcolsep}{2.6pt}
\begin{tabularx}{\linewidth}{@{}>{\raggedright\arraybackslash}p{2.75cm} *{8}{R}@{}}
\headrow \thead{€} & \thead{0} & \thead{1} & \thead{2} & \thead{3} & \thead{4} & \thead{5} & \thead{6} & \thead{7}\\
Ingresos & 0 & 1.500 & 5.000 & 30.000 & 50.000 & 35.000 & 30.000 & 0\\
\rowcolor{mist}$-$ Coste ventas & 0 & $-837$ & $-7.729$ & $-19.552$ & $-29.345$ & $-21.492$ & $-11.830$ & 0\\
$-$ Otros costes & $-4.000$ & $-2.200$ & $-1.200$ & $-1.300$ & $-1.500$ & $-1.600$ & $-1.700$ & 0\\
\rowcolor{mist}$-$ Amortización & 0 & $-2.250$ & $-2.250$ & $-2.250$ & $-2.250$ & $-2.250$ & $-2.250$ & 0\\
$+$ Plusvalía & & & & & & & & 500\\
\totrow \textbf{BAI} & $-4.000$ & $-3.787$ & $-6.179$ & 6.898 & 16.905 & 9.658 & 14.220 & 500\\
Impuesto 35\,\% & $-1.400$ & $-1.325,45$ & $-2.162,65$ & 2.414,30 & 5.916,75 & 3.380,30 & 4.977 & 175\\
\totrow \textbf{BDI} & $-2.600$ & $-2.461,55$ & $-4.016,35$ & 4.483,70 & 10.988,25 & 6.277,70 & 9.243 & 325\\
\end{tabularx}
\endgroup

\paso{b) Inversión inicial, flujos y valor residual}
\[
C_0\di=-[15.000+4.000\cdot(1-0,35)]=-(15.000+2.600)=\mathbf{-17.600}
\]

\begingroup\scriptsize\setlength{\tabcolsep}{2.6pt}
\begin{tabularx}{\linewidth}{@{}>{\raggedright\arraybackslash}p{2.75cm} *{7}{R}@{}}
\headrow \thead{€} & \thead{1} & \thead{2} & \thead{3} & \thead{4} & \thead{5} & \thead{6} & \thead{7}\\
$Q_i\ai$ & $-1.537$ & $-3.929$ & 9.148 & 19.155 & 11.908 & 16.470 & 0\\
\rowcolor{mist}$Q_i\ai\cdot0,65$ & $-999,05$ & $-2.553,85$ & 5.946,20 & 12.450,75 & 7.740,20 & 10.705,50 & 0\\
$+\,CAF\cdot0,35$ & 787,50 & 787,50 & 787,50 & 787,50 & 787,50 & 787,50 & 0\\
\totrow $Q_i\di$ renta & $-211,55$ & $-1.766,35$ & 6.733,70 & 13.238,25 & 8.527,70 & 11.493 & 0\\
$FM_i$ & 1.000 & 500 & 300 & 1.500 & 1.500 & 1.000 & 500\\
\rowcolor{mist}$-\Delta FM_i$ & $-1.000$ & $+500$ & $+200$ & $-1.200$ & 0 & $+500$ & $+500$\\
\totrow $Q_i\di$ tesorería & $-1.211,55$ & $-1.266,35$ & 6.933,70 & 12.038,25 & 8.527,70 & 11.993 & 500\\
\end{tabularx}
\endgroup

\smallskip
Valor residual en el periodo 7: activo fijo $VR_7\di=2.000-500\cdot0,35=1.825$; el FM que queda (500) se recupera por
su valor contable. Flujo total del periodo 7: $500+1.825+500=\mathbf{2.825}$.

\begin{trampa}[El FM del periodo 7]
En el periodo 7 el FM baja de 1.000 a 500 ($+500$ en tesorería) y los 500 restantes se recuperan como valor residual:
en total vuelven los 1.000 que había al cierre del periodo 6. No los cuentes dos veces ni te dejes la mitad.
\end{trampa}

\paso{c) VAN ($r=15\,\%$)}
\[
VAN=-17.600-\frac{1.211,55}{1,15}-\frac{1.266,35}{1,15^2}+\frac{6.933,70}{1,15^3}+\frac{12.038,25}{1,15^4}
+\frac{8.527,70}{1,15^5}+\frac{11.993}{1,15^6}+\frac{2.825}{1,15^7}=\mathbf{2.317,57}
\]

\begin{solucion}
$C_0\di=-17.600$; \ VAN $=2.317,57$\eur $>0$: el proyecto crea valor.
\end{solucion}

% ---------------------------------------------------------------------
\problema{18}{Estados financieros proyectados}{Problema · inversiones sucesivas leídas del balance}

\begin{enunciado}{18}
Determinada empresa, que se constituye a inicios de 2018, presenta los siguientes estados financieros proyectados, que
se corresponden con el proyecto de inversión a desarrollar:

\begin{center}\small
\begin{tabularx}{0.76\linewidth}{@{}L QQQ@{}}
\headrow \thead{Balance resumido} & \thead{Inicio 2018} & \thead{Final 2018} & \thead{Final 2019}\\
Activo circulante operativo & 100 & 200 & 300\\
Activo fijo bruto & 200 & 469 & 1.170\\
($-$) Fondo acumulado de depreciación & 0 & $-20$ & $-67$\\
\rowcolor{mist} Total activo & 300 & 649 & 1.403\\
Pasivo circulante operativo & 100 & 100 & 100\\
Deuda a largo plazo & 100 & 50 & 0\\
Fondos propios & 100 & 499 & 1.303\\
\rowcolor{mist} Total pasivo y neto & 300 & 649 & 1.403\\
\end{tabularx}
\end{center}

Además: a) las ventas para el año 2018 y 2019 se estiman en 1000 y 2000 euros respectivamente; b) los costes operativos
para el año 2018 y 2019 se estiman en un 40\,\% de las ventas; c) la amortización técnica de todas las inversiones que
realiza y realice en el futuro la empresa serán lineales en 10 años, sistema admitido a efectos fiscales; d) el coste
financiero de la deuda a largo plazo es de una tasa anual del 8\,\%; e) el impuesto de sociedades es del 30\,\%; f) el
coste de capital que aplica la empresa a todos sus proyectos de inversión es del 10\,\%; g) no estime el valor residual
del proyecto de inversión.

Con esta información, se pide calcular (siguiendo el orden indicado, y dejando en blanco aquel/aquellos puntos a los
que no sepa dar una respuesta): 1. Los ingresos y costes operativos anuales. 2. Las inversiones a realizar cada año.
3. La cuota de amortización técnica anual. 4. El beneficio antes de intereses e impuestos (BAIT) anual. 5. Los intereses
a pagar cada periodo. 6. El beneficio bruto anual. 7. El impuesto anual y el beneficio neto anual. 8. Las variaciones
del fondo de maniobra anuales. 9. El flujo de caja operativo después de impuestos a inicios de 2018, a finales de 2018
y a finales de 2019. 10. El VAN del proyecto. \textbf{SOL: VAN = $-$104,32\eur}
\end{enunciado}

\paso{1--3. Ingresos, costes, inversiones y amortización}
\begin{itemize}
  \item 2018: ventas 1.000, costes 400. \ 2019: ventas 2.000, costes 800.
  \item Inversiones (variación del activo fijo bruto): 200 en $t=0$; $469-200=269$ en $t=1$; $1.170-469=701$ en $t=2$.
  \item Amortización: la de 200 aporta 20/año desde 2018; la de 269 aporta 26,90 desde 2019; la de 701 empieza en 2020.
        Total: \textbf{20} en 2018 y \textbf{46,90} en 2019 (acumulada 66,90 $\approx$ 67, como en el balance).
\end{itemize}

\paso{4--7. Cuenta de resultados}
Intereses al 8\,\% sobre la deuda al inicio de cada año: $100\cdot0,08=8$ en 2018 y $50\cdot0,08=4$ en 2019.

\begin{center}\small
\begin{tabularx}{0.66\linewidth}{@{}L QQ@{}}
\headrow \thead{€} & \thead{2018} & \thead{2019}\\
Ventas & 1.000,00 & 2.000,00\\
\rowcolor{mist}$-$ Costes operativos (40\,\%) & $-400,00$ & $-800,00$\\
$-$ Amortización & $-20,00$ & $-46,90$\\
\totrow \textbf{4. BAIT} & 580,00 & 1.153,10\\
$-$ 5. Intereses & $-8,00$ & $-4,00$\\
\totrow \textbf{6. BAT} & 572,00 & 1.149,10\\
$-$ Impuesto (30\,\%) & $-171,60$ & $-344,73$\\
\totrow \textbf{7. Beneficio neto} & 400,40 & 804,37\\
\end{tabularx}
\end{center}

\paso{8. Fondo de maniobra (AC operativo $-$ PC operativo)}
$FM$: 0 (inicio 2018), 100 (final 2018), 200 (final 2019) $\Rightarrow$ $\Delta FM_1=+100$, $\Delta FM_2=+100$.

\paso{9. Flujos de caja (sin intereses) y 10. VAN}
\begin{center}\small
\begin{tabularx}{0.82\linewidth}{@{}L QQQ@{}}
\headrow \thead{€} & \thead{Inicio 2018} & \thead{Final 2018} & \thead{Final 2019}\\
$Q_i\ai=$ ventas $-$ costes & -- & 600,00 & 1.200,00\\
\rowcolor{mist}$Q_i\ai\cdot0,70$ & -- & 420,00 & 840,00\\
$+\,CAF\cdot0,30$ & -- & 6,00 & 14,07\\
\totrow $Q_i\di$ renta & -- & 426,00 & 854,07\\
$-\Delta FM_i$ & 0 & $-100,00$ & $-100,00$\\
\totrow \textbf{9. $Q_i\di$ tesorería} & 0 & 326,00 & 754,07\\
$-$ Inversión en activo fijo & $-200,00$ & $-269,00$ & $-701,00$\\
\totrow \textbf{Flujo neto total} & $-200,00$ & 57,00 & 53,07\\
\end{tabularx}
\end{center}
\[
VAN=-200+\frac{57,00}{1,10}+\frac{53,07}{1,10^2}=-200+51,82+43,86=\mathbf{-104,32}
\]

\begin{solucion}
VAN $=-104,32$\eur: con los datos de estos dos años, el proyecto no se acometería.
\end{solucion}
\begin{trampa}[Los intereses no van al flujo]
Los intereses se calculan porque los pide la cuenta de resultados (puntos 5--7), pero \textbf{no} se restan en el flujo
de caja del proyecto. Ahí solo entran las operaciones, el FM y las inversiones.
\end{trampa}

% ---------------------------------------------------------------------
\problema{19}{Ampliación de capacidad con préstamo}{Problema · separar inversión y financiación}

\begin{enunciado}{19}
Una empresa se está planteando realizar una inversión para incrementar su capacidad de producción. Para ello, deberá
adquirir dos máquinas que costarían 50.000\eur cada una, ambas amortizables linealmente en cinco años (sistema admitido a
efectos fiscales) y sin valor residual alguno. La realización del proyecto requeriría alquilar una nave contigua a las
instalaciones de la empresa durante los cinco años de duración del mismo, lo que supondría un desembolso anual de
12.000\eur.

Se espera obtener unos ingresos anuales por ventas de 150.000\eur, y se estiman unos costes operativos de 60.000\eur/año,
cifra que no incluye la amortización técnica de los equipos. El incremento inicial del FM imputable al proyecto es de
12.000\eur. Posteriormente, al final de cada uno de sus cinco años de duración, se estima que el FM ascenderá al 5\,\% del
volumen anual de ingresos previstos. El FM se recuperará al final de la vida del proyecto por su valor contable.

Para la financiación, se tiene previsto contratar un préstamo bancario a 5 años, por un importe de 50.000\eur, con
amortización financiera al vencimiento, y a un 6\,\% de interés anual. El coste de oportunidad del capital es del
10,5\,\%, y la tasa del Impuesto de Sociedades es del 35\,\%. Con esta información, se pide:
\begin{opciones}
  \item Calcular, para cada uno de los cinco años de duración del proyecto, el BAIT, el BAT y el BN imputables al mismo.
  \item Calcular el desembolso inicial, los flujos netos de caja de explotación, el valor residual y el VAN del
        proyecto. \textbf{SOL: VAN = 112.587,82\eur}
\end{opciones}
\end{enunciado}

\paso{a) Resultados (iguales los cinco años)}
$CAF=100.000/5=20.000$; intereses $=50.000\cdot0,06=3.000$.
\begin{center}\small
\begin{tabularx}{0.6\linewidth}{@{}L Q@{}}
\headrow \thead{€/año (años 1 a 5)} & \thead{Importe}\\
Ventas & 150.000\\
\rowcolor{mist}$-$ Costes operativos & $-60.000$\\
$-$ Alquiler de la nave & $-12.000$\\
\rowcolor{mist}$-$ Amortización & $-20.000$\\
\totrow \textbf{BAIT} & 58.000\\
$-$ Intereses & $-3.000$\\
\totrow \textbf{BAT} & 55.000\\
$-$ Impuesto 35\,\% & $-19.250$\\
\totrow \textbf{BN} & 35.750\\
\end{tabularx}
\end{center}

\paso{b) Desembolso inicial}
$C_0\di=-100.000-12.000=\mathbf{-112.000}$. El préstamo no aparece: la financiación ya está recogida en $r$.

\paso{Flujos de explotación}
$Q_i\ai=150.000-60.000-12.000=78.000$; \ renta $=78.000\cdot0,65+20.000\cdot0,35=\mathbf{57.700}$.\\
FM: de 12.000 en $t=0$ pasa a $5\,\%\cdot150.000=7.500$ $\Rightarrow$ $\Delta FM_1=-4.500$ (se libera
y suma). Tesorería: año 1: $57.700+4.500=62.200$; años 2--5: 57.700.

\paso{Valor residual}
Máquinas: 0. FM recuperado por su valor contable: 7.500. Año 5: $57.700+7.500=65.200$.

\begin{center}
\flujos[2.3]{5}{0/n/{$-112.000$},1/p/{62.200},2/p/{57.700},3/p/{57.700},4/p/{57.700},5/p/{57.700\\$+$7.500}}
\end{center}

\[
VAN=-112.000+\frac{62.200}{1,105}+\frac{57.700}{1,105^2}+\frac{57.700}{1,105^3}+\frac{57.700}{1,105^4}
+\frac{65.200}{1,105^5}=\mathbf{112.587,82}
\]

\begin{solucion}
$C_0\di=-112.000$; \ $Q\di$: 62.200 y 57.700; \ $VR_5\di=7.500$; \ VAN $=112.587,82$\eur.
\end{solucion}

% ---------------------------------------------------------------------
\problema{20}{Tienda online con canibalización}{Problema · VAN incremental}

\begin{enunciado}{20}
Una empresa de calzado deportivo se está planteando la posibilidad de introducirse en el comercio electrónico para
ampliar sus canales de distribución, por lo que está analizando el lanzamiento de un proyecto piloto para vender on line
parte de sus productos durante 3 años. El coste de inversión en la plataforma asciende a 10.000 euros (que se
amortizaría contable y fiscalmente en 2 años), no existiendo valor residual alguno al finalizar los 3 años del
proyecto. Además, se realizaría un curso inicial de formación al personal, con un coste de 2.000 euros, y se ampliarían
las existencias iniciales en 2.500 euros. Estas existencias serían recuperadas al final de la vida del proyecto por el
70\,\% de su valor. El gasto de mantenimiento de la plataforma supondría 1.000 euros anuales.

La cifra actual de negocios de la empresa a través de sus tiendas físicas es de 100.000 euros al año, importe que se
prevé estable para los próximos 3 años. Por su parte, las ventas online se estiman en 20.000 euros para el primer año,
con crecimientos del 5\,\% en los dos años siguientes. No obstante, si se incorpora la tienda virtual, se prevé que parte
de las ventas que actualmente se realizan a través del canal físico se deriven al canal electrónico, de modo que las
ventas del primero caerían en un 10\,\% para cada uno de los tres años. El resto de gastos de explotación no se verían
alterados.

Si el tipo del Impuesto de Sociedades es del 25\,\% (la empresa obtiene beneficios con su negocio actual, por el que tiene
previsto tributar en los próximos tres años), la financiación del proyecto es íntegramente con fondos propios y la tasa
de actualización adecuada para el proyecto es del 5\,\%, se pide:
\begin{enumerate}
  \item Cuenta de pérdidas y ganancias incremental para los ejercicios 0, 1, 2 y 3.
  \item Cálculo del VAN incremental, indicando la decisión que debería adoptar la empresa.
        \textbf{SOL: $\Delta$VAN = 10.388,30\eur}
\end{enumerate}
\end{enunciado}

\paso{Ventas incrementales}
Online: 20.000; 21.000; 22.050. Canibalización: $-10\,\%\cdot100.000=-10.000$ cada año.
Incremental: \textbf{10.000; 11.000; 12.050}.

\paso{1. Pérdidas y ganancias incremental}
\begin{center}\small
\begin{tabularx}{0.86\linewidth}{@{}L QQQQ@{}}
\headrow \thead{€} & \thead{Año 0} & \thead{Año 1} & \thead{Año 2} & \thead{Año 3}\\
Ventas incrementales & 0 & 10.000 & 11.000 & 12.050\\
\rowcolor{mist}$-$ Curso de formación & $-2.000$ & 0 & 0 & 0\\
$-$ Mantenimiento & 0 & $-1.000$ & $-1.000$ & $-1.000$\\
\rowcolor{mist}$-$ Amortización plataforma & 0 & $-5.000$ & $-5.000$ & 0\\
$-$ Pérdida en existencias (30\,\%) & 0 & 0 & 0 & $-750$\\
\totrow \textbf{$\Delta$BAIT} & $-2.000$ & 4.000 & 5.000 & 10.300\\
$-$ Impuesto 25\,\% & $+500$ & $-1.000$ & $-1.250$ & $-2.575$\\
\totrow \textbf{$\Delta$BN} & $-1.500$ & 3.000 & 3.750 & 7.725\\
\end{tabularx}
\end{center}

\paso{2. Flujos incrementales}
\begin{itemize}
  \item $t=0$: $-[10.000+2.000\cdot(1-0,25)+2.500]=\mathbf{-14.000}$ (el curso es un gasto deducible).
  \item $t=1$: $(10.000-1.000)\cdot0,75+5.000\cdot0,25=\mathbf{8.000}$.
  \item $t=2$: $(11.000-1.000)\cdot0,75+5.000\cdot0,25=\mathbf{8.750}$.
  \item $t=3$: $(12.050-1.000)\cdot0,75=8.287,50$, más el valor residual de las existencias:
        $1.750+750\cdot0,25=1.937,50$ $\Rightarrow$ \textbf{10.225}.
\end{itemize}

\begin{center}
\flujos[2.6]{3}{0/n/{$-14.000$},1/p/{8.000},2/p/{8.750},3/p/{8.287,50\\$+$1.937,50}}
\end{center}
\[
\Delta VAN=-14.000+\frac{8.000}{1,05}+\frac{8.750}{1,05^2}+\frac{10.225}{1,05^3}=\mathbf{10.388,30}
\]

\begin{solucion}
$\Delta VAN=10.388,30$\eur $>0$: conviene lanzar la tienda online, aun descontando las ventas que quita a las tiendas.
\end{solucion}

% ---------------------------------------------------------------------
\problema{21}{Sonora, S.A.: ampliación}{Problema · solo cuentan los flujos incrementales}

\begin{enunciado}{21}
La empresa Sonora, S.A. presenta en el momento actual el siguiente balance de situación:

\begin{center}\small
\begin{tabularx}{0.86\linewidth}{@{}L Q L Q@{}}
\headrow \thead{Activo ($t=0$)} & \thead{€} & \thead{Pasivo ($t=0$)} & \thead{€}\\
Activos financieros & 20.000.000 & Capital & 40.000.000\\
Activos explotación & 80.000.000 & Crédito a l/p & 30.000.000\\
 & & Reservas & 30.000.000\\
\rowcolor{mist}\textbf{Activo} & 100.000.000 & \textbf{Pasivo} & 100.000.000\\
\end{tabularx}
\end{center}

\begin{itemize}
  \item Los activos financieros del balance se refieren a una inversión que, por el momento, no piensa liquidarse, y que
        produce un 22\,\% de rentabilidad bruta anual postpagable.
  \item Los activos de explotación producen unos flujos de caja de explotación anuales de 20 millones de\eur, y se vienen
        amortizando, contable y fiscalmente, en 10 millones de\eur/año. Transcurridos tres años pueden venderse por 50
        millones de euros.
\end{itemize}
Dado el éxito empresarial alcanzado, la entidad se plantea la posibilidad de llevar a cabo una ampliación de su
actividad productiva en los siguientes términos:
\begin{itemize}
  \item Compra de equipos por 20 millones de\eur, a pagar al contado, amortizables linealmente en dos años. Esta cuota
        contable está permitida a efectos fiscales.
  \item Gastos de transporte de los equipos por una cuantía de 2 millones de\eur.
  \item La nueva inversión va a generar un volumen de ventas de 12 millones de\eur/año y unos gastos de explotación de 2
        millones\eur/año (sin amortización técnica).
  \item La duración prevista para esta inversión es de tres años, a final de los cuales se podrán vender los equipos
        nuevos por 6 millones de\eur.
  \item Incremento inicial del fondo de maniobra de 200.000\eur (en adelante no se producirá ningún incremento o reducción
        adicional en el fondo de maniobra). Este capital circulante se liquidaría al final de los tres años por su valor
        neto contable.
  \item La tasa impositiva es del 35\,\%.
  \item El coste de oportunidad del capital para el nuevo proyecto es del 4,73\,\%.
\end{itemize}
Con la información disponible, determine la conveniencia, o no, de llevar a cabo la ampliación.
\end{enunciado}

\begin{clave}[Primero, descartar lo irrelevante]
Los activos financieros y los de explotación actuales siguen igual se amplíe o no: sus flujos son comunes a las dos
alternativas y se cancelan. \textbf{Solo se valora la ampliación.}
\end{clave}

\paso{Inversión inicial incremental (criterio de la cátedra)}
$\Delta A_0=-(20.000.000+2.000.000+200.000)=\mathbf{-22.200.000}$

\paso{Flujos de explotación}
$CAF=20.000.000/2=10.000.000$ en los años 1 y 2; en el año 3 ya está amortizado.
\[
\Delta Q_{1,2}=10.000.000\cdot0,65+10.000.000\cdot0,35=\mathbf{10.350.000}\qquad
\Delta Q_{3}=10.000.000\cdot0,65=\mathbf{6.500.000}
\]

\paso{Valor residual en el año 3}
Equipos con valor contable 0 vendidos por 6.000.000: plusvalía íntegra. FM recuperado por su valor contable.
\[
\Delta VR_3=\underbrace{6.000.000-6.000.000\cdot0,35}_{3.900.000}+\underbrace{200.000}_{FM}=\mathbf{4.100.000}
\]

\begin{center}
\flujos[2.6]{3}{0/n/{$-22.200.000$},1/p/{10.350.000},2/p/{10.350.000},3/p/{6.500.000\\$+$4.100.000}}
\end{center}
\[
\Delta VAN=-22.200.000+\frac{10.350.000}{1,0473}+\frac{10.350.000}{1,0473^2}+\frac{10.600.000}{1,0473^3}
=\mathbf{6.346.457,61}
\]

\begin{solucion}
$\Delta VAN=6.346.457,61$\eur $>0$: conviene llevar a cabo la ampliación.
\end{solucion}
\begin{trampa}[Gastos de transporte]
La solución oficial suma los 2 M€ de transporte al desembolso inicial y no los amortiza ni los deduce. Si se trataran
como gasto deducible en $t=0$, el desembolso sería $-21.500.000$ y el $\Delta$VAN subiría a 7.046.457,61. En el examen,
sigue el criterio de la cátedra.
\end{trampa}

% ---------------------------------------------------------------------
\problema{22}{Sr.\ Tired: continuar o cesar}{Problema · dos alternativas excluyentes y riesgo distinto}

\begin{enunciado}{22}
El Sr.~Tired no confía demasiado en las expectativas del gobierno respecto de la incipiente recuperación del país, por
lo que se está planteando si continuar con su negocio de venta de libros durante dos ejercicios más, al final de los
cuales en cualquier caso se retirará, o bien cesar la actividad en el momento presente ante su previsión de
continuación en la tendencia de bajada en las ventas y resultados que viene acusando. Su balance en euros a 31/12/2013
(fecha que tomaremos como momento 0) es el siguiente:

\begin{center}\small
\begin{tabularx}{0.92\linewidth}{@{}L Q L Q@{}}
\headrow \thead{Activo} & \thead{€} & \thead{Pasivo} & \thead{€}\\
Inmovilizado material (local) & 500.000 & Capital social & 3.000\\
Mobiliario y equipamiento & 100.000 & Reservas & 1.197.000\\
Existencias & 400.000 & Resultados negativos ej.\ anteriores & (200.000)\\
\rowcolor{mist}\textbf{Total} & 1.000.000 & \textbf{Total} & 1.000.000\\
\end{tabularx}
\end{center}

En caso de continuar con la actividad, los datos que estima serán los siguientes:
\begin{itemize}
  \item La cifra de negocios prevista es de 600.000 euros en 2014 y un 20\,\% menos en 2015.
  \item Los costes de explotación supondrán un 75\,\% de las ventas en cada uno de los años, al que habrá que añadir el
        coste de despido de los trabajadores al final del 2015, que se estima ascenderá a 200.000 euros.
  \item La empresa desarrolla su actividad en un local propiedad de la sociedad que ya se encuentra totalmente
        amortizado a efectos contables y fiscales.
  \item El mobiliario y equipamiento se amortizarán de forma lineal en dos años, siendo esta amortización admitida como
        gasto deducible por la administración tributaria. Transcurridos los dos años, no se prevé obtener ningún valor
        por estos activos.
  \item Se prevé una reducción del fondo de maniobra de 150.000 euros en el primer año y la liquidación del mismo por
        su valor contable al final del 2015.
\end{itemize}
En caso de cesar en la actividad en el momento actual, el local (sin mobiliario ni equipamiento) sería alquilado,
durante un periodo de dos años, a un tercero muy solvente que ya ha manifestado su interés y con el que se ha negociado
el pago de una renta anual, avalado personalmente por el arrendatario, del 5\,\% del valor neto contable del inmueble.
El personal sería despedido de forma inmediata, lo que supondría un coste de indemnizaciones 175.271,55 euros. El resto
de activos serían liquidados de acuerdo a la siguiente estimación:
\begin{itemize}
  \item Las existencias se podrían vender hoy por un precio aproximado del 110\,\% de su valor contable.
  \item El mobiliario y equipamiento se venderían por 75.000 euros.
\end{itemize}
Si la tasa impositiva es del 30\,\%, se pide:
\begin{enumerate}
  \item Indique las alternativas que se está planteando el Sr.~Tired.
  \item Calcule los flujos netos de caja después de impuestos de las dos alternativas.
  \item Si la tasa de actualización fuese del 10\,\% en cualquiera de las dos alternativas, calcule el valor
        actualizado neto (VAN) incremental. \textbf{SOL: $\Delta$VAN = 0\eur}
  \item Comente los resultados obtenidos y explique razonadamente si le parece adecuado usar la misma tasa de descuento
        para valorar las dos alternativas.
\end{enumerate}
\end{enunciado}

\paso{1. Alternativas}
\textbf{I}: continuar dos años más (2014 y 2015) y retirarse. \ \textbf{II}: cesar hoy, liquidar mobiliario y existencias,
despedir al personal y alquilar el local dos años.

\paso{2a. Alternativa I: continuar}
$CAF=100.000/2=50.000$ $\Rightarrow$ ahorro fiscal 15.000/año. FM (existencias): $400.000\to250.000\to0$.
\begin{center}\small
\begin{tabularx}{0.84\linewidth}{@{}L QQQ@{}}
\headrow \thead{Alternativa I (€)} & \thead{2013} & \thead{2014} & \thead{2015}\\
Ventas & & 600.000 & 480.000\\
\rowcolor{mist}$-$ Costes de explotación (75\,\%) & & $-450.000$ & $-360.000$\\
$-$ Despidos & & & $-200.000$\\
\rowcolor{mist}$Q_i\ai$ & & 150.000 & $-80.000$\\
$Q_i\ai\cdot0,70$ & & 105.000 & $-56.000$\\
\rowcolor{mist}$+\,CAF\cdot0,30$ & & 15.000 & 15.000\\
\totrow $Q_i\di$ renta & & 120.000 & $-41.000$\\
$FM_t$ & 400.000 & 250.000 & 0\\
\rowcolor{mist}$-\Delta FM_t$ & & $+150.000$ & $+250.000$\\
\totrow \textbf{Flujo total I} & 0 & 270.000 & 209.000\\
\end{tabularx}
\end{center}

\paso{2b. Alternativa II: cesar y alquilar}
\begin{itemize}
  \item Despidos hoy: $-175.271,55\cdot0,70=-122.690,09$.
  \item Existencias al 110\,\%: $440.000-40.000\cdot0,30=428.000$ (plusvalía de 40.000).
  \item Mobiliario por 75.000 (valor contable 100.000): $75.000+25.000\cdot0,30=82.500$ (minusvalía de 25.000).
  \item Total en $t=0$: $-122.690,09+428.000+82.500=\mathbf{387.809,91}$.
  \item Alquiler en 2014 y 2015: $5\,\%\cdot500.000=25.000$ $\Rightarrow$ $25.000\cdot0,70=\mathbf{17.500}$/año.
\end{itemize}

\begin{center}
\begin{minipage}{0.48\linewidth}\centering
{\sffamily\scriptsize\bfseries\color{navy}Alternativa I: continuar}\par\smallskip
\flujos[1.7]{2}{1/p/{270.000},2/p/{209.000}}
\end{minipage}\hfill
\begin{minipage}{0.48\linewidth}\centering
{\sffamily\scriptsize\bfseries\color{navy}Alternativa II: cesar y alquilar}\par\smallskip
\flujos[1.7]{2}{0/p/{387.809,91},1/p/{17.500},2/p/{17.500}}
\end{minipage}
\end{center}

\paso{3. VAN ($r=10\,\%$)}
\[
VAN_I=\frac{270.000}{1,10}+\frac{209.000}{1,10^2}=418.181,82
\qquad
VAN_{II}=387.809,91+\frac{17.500}{1,10}+\frac{17.500}{1,10^2}=418.181,82
\]
\[
\Delta VAN_{II/I}=\mathbf{0}
\]

\paso{4. Comentario}
Con la misma tasa, las dos alternativas valen lo mismo. Pero \textbf{no es adecuado} usar la misma $r$: vender libros
con ventas en caída es mucho más arriesgado que cobrar un alquiler de un inquilino muy solvente que además lo avala
personalmente. La alternativa II merece una tasa menor ($r_{II}<r_I$), con lo que $VAN_{II}>VAN_I$ y
\textbf{convendría cesar y alquilar}.

\begin{solucion}
$\Delta VAN=0$ con $r=10\,\%$; corrigiendo por riesgo, es preferible cesar y alquilar el local.
\end{solucion}
\begin{trampa}[La renta del local]
El enunciado dice que el local está totalmente amortizado, pero el balance lo recoge por 500.000. La solución oficial
toma esos 500.000 como valor neto contable: renta de 25.000 al año.
\end{trampa}

% ---------------------------------------------------------------------
\problema{23}{Para qué sirve el CAE}{Test · teoría del coste anual equivalente}

\begin{enunciado}{23}
Indique cuál de las siguientes afirmaciones sobre el cálculo del coste anual equivalente es cierta:
\begin{opciones}
  \item No resulta conveniente calcular el coste anual equivalente en anualidades reales, porque puede proporcionar
        resultados incorrectos cuando las tasas de inflación son elevadas.
  \item Este cálculo sólo es apropiado cuando estamos comparando proyectos de inversión mutuamente excluyentes, pero de
        idéntica duración o estructura temporal.
  \item El criterio del VAN es más adecuado que el del coste anual equivalente en cualquier circunstancia. Si existe
        discrepancia entre ambos cálculos. Siempre emplearemos el primer método de valoración.
  \item Ninguna de las anteriores es correcta.
\end{opciones}
\end{enunciado}

\begin{itemize}
  \item \textbf{a) es falsa}: con inflación, lo correcto es precisamente calcular el CAE en términos \textbf{reales}
        ($r_{real}=\frac{1+r_{nom}}{1+\pi}-1$); en nominal la cuota no sería constante.
  \item \textbf{b) es falsa}: el CAE existe para comparar equipos de \textbf{distinta} duración. Si duran lo mismo,
        basta con el VA de los costes.
  \item \textbf{c) es falsa}: con equipos de distinta vida que se reponen, comparar el VAN de un solo ciclo lleva a
        error; hay que usar el CAE.
\end{itemize}

\begin{solucion}
\textbf{d)}: ninguna es correcta.
\end{solucion}

% ---------------------------------------------------------------------
\problema{24}{Quickpizza, S.L.: horno A u horno B}{Problema · CAE con vidas distintas}

\begin{enunciado}{24}
Quickpizza, S.L. se está planteando la sustitución de su horno actual por otro nuevo. El gerente ha recibido las dos
ofertas siguientes:
\begin{itemize}
  \item Horno A, con un coste inicial de 17.000\eur, costes operativos anuales de 3.000\eur, y vida útil de 4 años.
  \item Horno B, con un coste inicial del 12.000\eur, costes operativos anuales de 4.000\eur, y vida útil de 3 años.
\end{itemize}
Si el coste de oportunidad del capital es del 12\,\%, ¿qué horno debería elegir la empresa? Justifique numéricamente
su respuesta. \textbf{SOL: CAE$_A$ = 8.596,99\eur; CAE$_B$ = 8.996,19\eur}
\end{enunciado}

Vidas distintas y equipos que se reponen: se compara el \textbf{coste anual equivalente}.

\begin{center}\small
\begin{tabularx}{0.9\linewidth}{@{}l L L@{}}
\headrow & \thead{Horno A ($n=4$)} & \thead{Horno B ($n=3$)}\\
Factor $\an{n}{0{,}12}$ & 3,037349 & 2,401831\\
\rowcolor{mist} VA de costes & $17.000+3.000\cdot3,037349=26.112,05$ & $12.000+4.000\cdot2,401831=21.607,33$\\
CAE $=$ VA$/\an{n}{r}$ & $26.112,05/3,037349=\mathbf{8.596,99}$ & $21.607,33/2,401831=\mathbf{8.996,19}$\\
\end{tabularx}
\end{center}

\begin{center}
\begin{tikzpicture}[font=\sffamily\scriptsize,x=0.82cm,y=0.6cm]
  \foreach \x/\l in {0/0,2/2.000,4/4.000,6/6.000,8/8.000,10/10.000}{
    \draw[rule] (\x,-0.25) -- (\x,1.75); \node[slate,below] at (\x,-0.25) {\l};}
  \fill[okgreen] (0,0.95) rectangle (8.597,1.55);
  \fill[slate] (0,0.15) rectangle (8.996,0.75);
  \node[anchor=east,navy] at (-0.15,1.25) {Horno A};
  \node[anchor=east,navy] at (-0.15,0.45) {Horno B};
  \node[anchor=west,fill=white,inner sep=1.5pt,font=\sffamily\scriptsize\bfseries,okgreen!70!black] at (8.65,1.25) {8.596,99};
  \node[anchor=west,fill=white,inner sep=1.5pt,font=\sffamily\scriptsize\bfseries,slate] at (9.05,0.45) {8.996,19};
  \draw[slate] (0,-0.25) -- (0,1.75);
  \node[slate] at (5,-1.15) {coste anual equivalente (€/año)};
\end{tikzpicture}
\end{center}

\begin{solucion}
Elegir el \textbf{horno A}: su CAE es menor (8.596,99 < 8.996,19; ahorra 399,20\eur al año).
\end{solucion}
\begin{trampa}[No compares el VA de un solo ciclo]
B tiene menos VA de costes (21.607 < 26.112) solo porque dura un año menos. Comparar esos VA llevaría a elegir B, que es
el error que quiere provocar el problema.
\end{trampa}

% ---------------------------------------------------------------------
\problema{25}{Copiadoras con demanda estacional}{Test · cuántas máquinas sustituir}

\begin{enunciado}{25}
Una empresa de reprografía, con demanda estacional, que actualmente tiene dos copiadoras en funcionamiento, se está
planteando la posibilidad de sustituir una o dos de estas máquinas por otras nuevas. Las características de las
copiadoras se detallan en la tabla siguiente:

\begin{center}\small
\begin{tabularx}{0.82\linewidth}{@{}L QQ@{}}
\headrow \thead{Concepto} & \thead{Copiadoras actuales} & \thead{Copiadoras nuevas}\\
Capacidad anual por máquina & 100.000 copias & 100.000 copias\\
\rowcolor{mist} Coste adquisición/máquina nueva & -- & 1.000.000\eur\\
Coste operativo por copia & 3\eur & 1,5\eur\\
\rowcolor{mist} Valor residual por máquina & -- & --\\
\end{tabularx}
\end{center}

Si todas las copiadoras tienen una vida indefinida, la demanda es de 150.000 copias para el período octubre-junio, y de
12.500 copias en verano (julio-septiembre), y el coste de capital es del 10\,\%, indique qué opción recomendaría a esta
empresa:
\begin{opciones}
  \item No sustituir ninguna copiadora.
  \item Sustituir ambas copiadoras.
  \item Sustituir una de las dos, lo que permitiría obtener un ahorro de costes de 450.000\eur si se utiliza la copiadora
        nueva a pleno rendimiento, y la antigua se destina a absorber el exceso en la época de mayor demanda.
  \item Ninguna de las opciones anteriores es correcta.
\end{opciones}
\end{enunciado}

Método de las diapositivas del Tema 4 (52--55): demanda anual $150.000+12.500=162.500$ copias; vida indefinida, así que
los costes anuales se actualizan como perpetuidad ($\div0,10$).

\begin{center}\small
\begin{tabularx}{\linewidth}{@{}L >{\raggedleft\arraybackslash}p{5.4cm} Q@{}}
\headrow \thead{Alternativa} & \thead{VA de costes} & \thead{Ahorro vs.\ I}\\
I. No sustituir (162.500 copias a 3\eur) & $\frac{487.500}{0,10}=4.875.000$ & --\\
\rowcolor{mist} II. Sustituir las dos (162.500 copias a 1,5\eur) & $2.000.000+\frac{243.750}{0,10}=4.437.500$ & 437.500\\
III. Sustituir una: nueva a pleno (100.000) $+$ vieja el resto (62.500) & $1.000.000+\frac{150.000}{0,10}+\frac{187.500}{0,10}=4.375.000$ & \textbf{500.000}\\
\end{tabularx}
\end{center}

\begin{center}
\begin{tikzpicture}[font=\sffamily\scriptsize,x=1.5cm,y=0.6cm]
  \foreach \x/\l in {0/0,1/1,2/2,3/3,4/4,5/5}{
    \draw[rule] (\x,-0.25) -- (\x,2.55); \node[slate,below] at (\x,-0.25) {\l};}
  \fill[slate] (0,1.75) rectangle (4.875,2.35);
  \fill[steel] (0,0.95) rectangle (4.4375,1.55);
  \fill[okgreen] (0,0.15) rectangle (4.375,0.75);
  \node[anchor=east,navy] at (-0.1,2.05) {I. No sustituir};
  \node[anchor=east,navy] at (-0.1,1.25) {II. Sustituir las dos};
  \node[anchor=east,navy] at (-0.1,0.45) {III. Sustituir una};
  \node[anchor=west,fill=white,inner sep=1.5pt,font=\sffamily\scriptsize\bfseries,slate] at (4.9,2.05) {4.875.000};
  \node[anchor=west,fill=white,inner sep=1.5pt,font=\sffamily\scriptsize\bfseries,steel] at (4.46,1.25) {4.437.500};
  \node[anchor=west,fill=white,inner sep=1.5pt,font=\sffamily\scriptsize\bfseries,okgreen!70!black] at (4.4,0.45) {4.375.000 (mínimo)};
  \draw[slate] (0,-0.25) -- (0,2.55);
  \node[slate] at (2.5,-1.15) {VA de los costes (millones de €)};
\end{tikzpicture}
\end{center}

Lo mejor es sustituir \textbf{una} (III), pero el ahorro es de \textbf{500.000\eur}, no de 450.000. Por eso c) no es
correcta tal y como está escrita.

\begin{solucion}
\textbf{d)}: conviene sustituir una sola copiadora, con un ahorro de 500.000\eur (62.500 frente a sustituir las dos).
\end{solucion}
\begin{trampa}[No hay solución oficial publicada de P25]
Esta respuesta aplica el método de las diapositivas, que reparte la capacidad \textbf{por año}. Si se exigiera
capacidad \textbf{por estación} (75.000 copias por máquina de octubre a junio y 25.000 en verano), la nueva solo podría
hacer $75.000+12.500=87.500$ copias, el ahorro de sustituir una bajaría a 312.500 y saldría mejor sustituir las dos
(437.500). Ninguna lectura da 450.000, así que c) es falsa en ambos casos. Si en clase se corrigió de otra forma,
manda el criterio del profesor.
\end{trampa}

\end{document}
